\documentclass[10pt,a4paper]{amsart}
\usepackage{amsmath,amssymb,amsthm,graphicx}
\usepackage[T1]{fontenc}
\usepackage[utf8]{inputenc}
\usepackage[english]{babel}
\usepackage{amsmath,amssymb,amsthm,mathtools,enumitem,microtype}
\usepackage[left=2.5cm,right=2.5cm,top=3cm,bottom=3cm]{geometry}
\usepackage{xcolor}
\usepackage[colorlinks=true,linkcolor=blue!50!black,citecolor=blue!50!black,urlcolor=blue!50!black]{hyperref}
\hypersetup{pdftitle={Amenable C-sequences and ineffability in L},
  pdfauthor={Gabriel Fernandes}}
\newtheorem{theorem}{Theorem}[section]
\newtheorem{lemma}[theorem]{Lemma}
\newtheorem{proposition}[theorem]{Proposition}
\newtheorem{corollary}[theorem]{Corollary}
\theoremstyle{definition}
\newtheorem{definition}[theorem]{Definition}
\theoremstyle{remark}
\newtheorem{remark}[theorem]{Remark}
\newcommand{\A}{\mathcal A}
\newcommand{\D}{\mathbb D}
\newcommand{\PP}{\mathbb P}
\newcommand{\RR}{\mathbb R}
\newcommand{\EE}{\mathbb E}
\newcommand{\ZFC}{\mathsf{ZFC}}
\newcommand{\ZFCc}{\mathsf{ZFC}^{-}_{\mathrm{Coll}}}
\DeclareMathOperator{\crit}{crit}
\DeclareMathOperator{\dom}{dom}
\DeclareMathOperator{\tc}{tc}
\DeclareMathOperator{\supp}{supp}
\DeclareMathOperator{\Ord}{Ord}
\DeclareMathOperator{\Inacc}{Inacc}
\DeclareMathOperator{\Comp}{Comp}
\DeclareMathOperator{\acc}{acc}

\title{Consistency of no amenable C-sequence with higher-order indescribability and non-ineffability in L}
\author{Gabriel Fernandes}
\date{arxiv:2610.09086v1; CC BY 4.0}
\begin{document}
\maketitle
\noindent\textit{Editorial note.} Complete original mathematical body and embedded bibliography retained. Author takes responsibility for mathematical correctness; no proof repair, additional assumptions or independent certification.\par\medskip
\noindent\textit{Editorial note.} Original declaration2714--2722 retained: Codex Sol6.1 assisted literature, strategy exploration, mathematical argument development and proof checks, organization and preparation; author takes full responsibility. No explicit core-proof AI-discovery target is asserted, and no AI-free/pure-human claim is added.\par\medskip
\noindent\textit{Editorial note.} Full preservation appendix explicitly included in the relevant proof bundles; the higher-order Pi\textasciicircum{}2\_1 convention, local lottery preparation and continuum assumptions remain unchanged.\par\medskip
\noindent\textit{Editorial note.} Three native statements are collected once each; direct capture/iteration corollaries and the alternative stronger-cardinal consistency route are not separately counted.\par\medskip
\noindent\textit{Allowed standard prerequisites (editorial).}\begin{itemize}
\item Higher-order indescribability, weak compactness, generalized Hechler forcing and elementary embedding lifts. Native Pi\textasciicircum{}2\_1 convention is explicitly distinguished from Pi\textasciicircum{}1\_2; all C-sequence amenability definitions, support/localization bounds, canonical presentation, target-model agreement and local lottery preparation assumptions retained.. Complete original body and bibliography, including the full nonadjacent preservation appendix2432--2673; no preservation assumption replaced by a stronger one.
\end{itemize}\medskip
\section{Introduction}\label{sec:introduction}

Conjecture 4.13 of \cite{IR} asks whether a regular uncountable
cardinal which carries no amenable \(C\)-sequence must be
ineffable in \(L\). We give a relative consistency
counterexample in which the cardinal remains
\(\Pi^2_1\)-indescribable. 

\begin{definition}[Notation and \(C\)-sequences]\label{def:c-sequence}
Throughout the article, inaccessibility means strong
inaccessibility, and \(\Inacc^M(\kappa)\) denotes the set of
inaccessible cardinals below \(\kappa\) as computed in the
specified universe \(M\). Without a superscript it is computed
in the current universe. For \(X\subseteq\eta\), where
\(\eta\) is a limit ordinal, write
\(\acc(X)=\{\delta<\eta:0<\delta=\sup(X\cap\delta)\}\),
with the ambient ordinal \(\eta\) understood from context.
For a regular uncountable cardinal \(\kappa\), a
\emph{\(C\)-sequence on \(S\subseteq\kappa\)} is a sequence
\(\vec C=\langle C_\beta:\beta\in S\rangle\) such that
\(C_\beta\) is club in \(\beta\) at every nonzero limit
index. At other indices in its domain we require
\(C_0=\varnothing\) and \(C_{\alpha+1}=\{\alpha\}\).
\end{definition}

\begin{definition}[Amenability and capture]
\label{def:amenability}
Let \(\kappa\) be regular uncountable and let \(\vec C\) be a
\(C\)-sequence on \(S\subseteq\kappa\). For a club
\(D\subseteq\kappa\), set
\[
 \operatorname{Cap}(D,\vec C)
   =\{\beta\in S:D\cap\beta\subseteq C_\beta\}.
\]
The sequence is \emph{amenable in \(\kappa\)} if this set is
nonstationary in \(\kappa\) for every club \(D\subseteq\kappa\).
We say that \(D\) \emph{captures} \(\vec C\) if
\(\operatorname{Cap}(D,\vec C)\) is stationary in \(\kappa\).
Following \cite[Definitions 4.1 and 4.3]{IR}, write
\[
 \A_\kappa
   =\{S\subseteq\kappa:S\text{ carries a \(C\)-sequence
                         amenable in \(\kappa\)}\}.
\]
\end{definition}

The notion of amenability was made explicit in \cite{BR};
we use the equivalent formulation recorded in
\cite[Definition 4.1 and Fact 4.2]{IR}.
The entries at zero and successor ordinals do not affect
amenability, since the nonzero limit ordinals form a club.
The family \(\A_\kappa\) is a \(\kappa\)-complete ideal
containing the nonstationary ideal. Restrictions preserve
amenability. For a union of fewer than \(\kappa\) sets carrying
amenable sequences, assign each index to the first set
containing it in a fixed enumeration, and use that set's
sequence there. For any club, the capture set of the resulting
sequence is contained in the union of the original capture
sets, a union of fewer than \(\kappa\) nonstationary sets.
A sequence on a nonstationary domain is automatically amenable.
In particular, \(\kappa\notin\A_\kappa\) says exactly that
every \(C\)-sequence on \(\kappa\) is captured by some club.


\begin{theorem}\label{thm:main}
If \(\ZFC+\) ``there is a \(\Pi^2_1\)-indescribable cardinal''
is consistent, then so is
\[
 \ZFC+\exists\kappa\,
 \bigl[\kappa\text{ is }\Pi^2_1\text{-indescribable}
       \ \land\ \kappa\notin\A_\kappa
       \ \land\ L\models\text{\(\kappa\) is not ineffable}\bigr].
\]
\end{theorem}


The converse consistency implication follows from the
indescribability asserted in the conclusion, so the two
theories in Theorem~\ref{thm:main} are equiconsistent.
Theorem~\ref{thm:main} also gives a negative answer to
\cite[Question~5.5]{BNZRad}, which asks whether the existence
of a cardinal carrying no amenable \(C\)-sequence is
equiconsistent with the existence of an ineffable cardinal.
Indeed, by \cite[Lemma~5.4]{HLN}, every ineffable cardinal
\(\delta\) is a limit of totally indescribable cardinals.
In particular, \(V_\delta\) satisfies \(\ZFC+\) ``there is
a \(\Pi^2_1\)-indescribable cardinal''. Thus the consistency
upper bound supplied by our construction is strictly below
an ineffable cardinal. We do not determine the optimal
consistency strength when the indescribability requirement
in the conclusion is omitted. The superscript in \(\Pi^2_1\)
is significant: we use higher-order indescribability,
as specified in Definition~\ref{def:indescribable},
rather than \(\Pi^1_2\)-indescribability.

Models with no amenable \(C\)-sequence were obtained by Radin
forcing \cite{BNZRad}, and later by overlapping-extender
forcing \cite{BNZApprox}. We obtain such a model by Hechler
forcing, while preserving \(\Pi^2_1\)-indescribability and
non-ineffability in \(L\). We prove stationary local
domination at weakly compact cardinals.

A related result in \cite{RYY} gives a model in which all
\(C\)-sequences are trivial in the sense of the
\(C\)-sequence number, while a \(\kappa\)-Souslin tree
exists. Triviality allows an initial segment of a cofinal
set to be covered by an entry chosen at an index that need
not equal the endpoint of that initial segment;
capture requires a club \(D\) such that
\(D\cap\beta\subseteq C_\beta\) for stationarily many
\(\beta\). Our model retains weak
compactness, and every \(C\)-sequence is captured in the
sense of Definition~\ref{def:amenability}.

We start in \(L\), fix its least
\(\Pi^2_1\)-indescribable cardinal \(\kappa\), and
show that \(\kappa\) is not ineffable there. We then
apply the local lottery preparation of \cite{Johnstone} and iterate
the full generalized Hechler forcing with support
\(<\kappa\) for \(\kappa^+\) stages. Each final
\(C\)-sequence appears in an intermediate extension;
the next Hechler stage captures it, and the closed tail
preserves the stationary capture set.

At height \(\kappa+1\), \cite[Corollary 38]{Johnstone} covers
orders of size at most \(\kappa\). Under \(2^\kappa=\kappa^+\),
Theorem~\ref{thm:preparation} extends this result to the full
Hechler iteration. Its proof in Appendix~\ref{sec:preservation}
uses a canonical presentation with conditions in \(H_{\kappa^+}\),
shows that the prepared universe and the embedding target
compute the same canonical iteration, and carries out two lifts.
For a \(\Pi^3_1\)-indescribable cardinal, the corresponding
local preparation has the larger size bound
\(\beth_{\kappa+1}=2^\kappa\), so Johnstone's corollary
applies directly to a dense presentation of the iteration.
Proposition~\ref{prop:pi3-preparation} and
Corollary~\ref{cor:pi3-consistency} give this alternative argument.

The use of a prepared cardinal followed by a long
Hechler iteration also appears in
\cite[Theorem 4.7]{Schuerz}, for a conclusion about
cardinal characteristics. There, indestructible strong
unfoldability is followed by a \(<\kappa\)-support
Hechler iteration of length \(\kappa^{++}\).
Here we use the weaker local hypothesis
\(\Pi^2_1\)-indescribability and length
\(\kappa^+\). The capture theorem ensures that every
\(C\)-sequence in the final extension is non-amenable.

Another use of Hechler domination at stationary sets of indices appears in
the proof of \cite[Lemma 20]{LL}.
In that construction, a classical Hechler real dominates
functions in \({}^\omega\omega\) associated with stationarily
many indices below \(\omega_1\); a subsequent forcing adds a
club through those indices. Our local domination theorem applies
to every ground-model family
\(\langle f_\beta:\beta\in\Inacc(\kappa)\rangle\),
without coherence assumptions on its entries.

The connection between Hechler forcing and club capture is
given by next-point functions. For inaccessible \(\beta\)
and \(\xi<\beta\), put
\(n_\beta(\xi)=\min(C_\beta\setminus(\xi+1))\).
If \(d[\beta]\subseteq\beta\) and
\(n_\beta\leq^*d\restriction\beta\) hold at stationarily
many inaccessible indices, the pressing-down lemma makes the domination
threshold uniform on a stationary subset. A single tail of the
club of limit closure points of \(d\) then captures \(\vec C\),
as shown in Corollary~\ref{cor:capture}. Ordinary eventual
domination of ground-model functions from \(\kappa\) to \(\kappa\)
does not by itself establish this local property: extending
\(n_\beta\) to domain \(\kappa\) may place the eventual
domination threshold beyond \(\beta\). The reflection argument
at a weakly compact cardinal supplies the stationary local
domination needed in Theorem~\ref{thm:domination}.

\section{Indescribability and the initial cardinal}\label{sec:initial}

\begin{definition}[Indescribability conventions]
\label{def:indescribable}
We use the classical notation \(\Pi^m_n\), in which the upper
index refers to quantifiers of order \(m+1\), with full
semantics in the ambient universe. Thus \(\Pi^1_1\) sentences
have a universal second-order block and a first-order matrix;
\(\Pi^2_1\) sentences permit a leading universal block of
third-order quantifiers. In particular, \(\Pi^2_1\) and
\(\Pi^1_2\) are different notions.
For inaccessible \(\kappa\), \(\Pi^{m+1}_1\)-indescribability
is equivalent to \((\kappa+m)\)-strong unfoldability for
\(1\leq m<\omega\), by \cite[Corollary 35]{Johnstone}.
Here \(\kappa+m\) uses ordinal addition. The two preservation
arguments below use heights \(\kappa+1\) and \(\kappa+2\),
respectively.
\end{definition}

We use weak compactness in its reflection form: if
\(R\subseteq V_\kappa\) and a \(\Pi^1_1\) sentence holds
in \((V_\kappa,\in,R)\), then it holds in
\((V_\beta,\in,R\cap V_\beta)\) for some \(\beta<\kappa\).
Every \(\Pi^2_1\)-indescribable cardinal, and hence every
\(\Pi^3_1\)-indescribable cardinal, is weakly compact.

\begin{proposition}\label{prop:least}
The least \(\Pi^2_1\)-indescribable cardinal, if it exists,
is neither subtle nor ineffable.
\end{proposition}
\begin{proof}
If the least such cardinal \(\kappa\) were subtle,
\cite[Lemma 5.4]{HLN} would give a totally indescribable
cardinal \(\gamma<\kappa\). In particular, \(\gamma\)
would be \(\Pi^2_1\)-indescribable, a contradiction.
Since ineffability implies subtlety, \(\kappa\) is not
ineffable.
\end{proof}

\begin{corollary}\label{cor:ground}
The consistency of a \(\Pi^2_1\)-indescribable cardinal
implies the consistency of a universe
\(W\models\ZFC+V=L\) containing such a cardinal.
\end{corollary}
\begin{proof}
The indescribability hypotheses used in \cite{Johnstone}
are, if consistent with \(\ZFC\), consistent with \(V=L\);
see \cite[Section 1, p.~1215]{Johnstone}.
\end{proof}

In that universe we take the least such \(\kappa\).
Proposition~\ref{prop:least} makes it non-ineffable in
\(L\), and GCH gives \(2^\kappa=\kappa^+\).

\section{Stationary capture by generalized Hechler forcing}
\label{sec:capture}

We prove local domination for a ground-model family of functions
indexed by the inaccessible cardinals below a weakly compact cardinal.
We first reflect predensity from \(\D_\kappa\) to \(\D_\beta\)
for an inaccessible \(\beta<\kappa\), then construct a condition
in \(\D_\kappa\) with stem length \(\beta\). Applied to the next-point
functions of a \(C\)-sequence,
this gives stationary capture.

\begin{definition}[Generalized Hechler forcing]\label{def:hechler}
For a regular uncountable cardinal \(\kappa\), let
\(\D_\kappa\) be generalized Hechler forcing in the presentation
of \cite[Definition~48]{BBFM}. Write its conditions as
\(p=(s_p,h_p)\), where \(s_p\in{}^{<\kappa}\kappa\),
\(h_p\in{}^\kappa\kappa\), and \(s_p\subseteq h_p\).
We call \(s_p\) the \emph{stem}, put \(\ell(p)=\dom(s_p)\),
and call \(h_p\) the \emph{side function}. Thus \(q\leq p\)
means \(s_p\subseteq s_q\) and \(h_p\leq h_q\) pointwise.
The definition makes sense for every regular uncountable
\(\kappa\); the hypothesis \(2^{<\kappa}=\kappa\) in the
reference is needed for the chain condition.
If \(\beta\leq\kappa\) is regular uncountable and
\(p\in\D_\kappa\) satisfies \(\ell(p)<\beta\) and
\(h_p[\beta]\subseteq\beta\), put
\[
 p\restriction\beta=(s_p,h_p\restriction\beta).
\]
Then \(p\restriction\beta\in\D_\beta\). We use this
notation only under the stated hypotheses.
\end{definition}

\begin{definition}[Domination and the generic function]
\label{def:eventual-domination}\label{def:generic-function}
For ordinal-valued functions \(f,g\) on a nonzero limit ordinal
\(\eta\), write \(f\leq^*g\) if
\[
 (\exists\gamma<\eta)(\forall\xi\in[\gamma,\eta))\,
 f(\xi)\leq g(\xi).
\]
If \(G\subseteq\D_\kappa\) is generic, its \emph{Hechler generic
function} is \(d_G=\bigcup_{p\in G}s_p\). In forcing statements,
\(d\) denotes the canonical name for this function.
\end{definition}

\begin{definition}[Predensity]\label{def:compatibility}
\label{def:predensity}
For a forcing order \(\mathbb P\), let
\(\Comp_{\mathbb P}(p,q)\) mean that \(p,q\in\mathbb P\) are
compatible. A family \(A\subseteq\mathbb P\) is \emph{predense
below \(p\)} if every \(q\leq p\) is compatible with a member of
\(A\); members of \(A\) need not themselves extend \(p\).
We omit the subscript on \(\Comp\) when the order is clear.
\end{definition}

\begin{remark}\label{rem:compatibility}
For \(p=(s,h),q=(t,k)\in\D_\kappa\),
\begin{align}
 \Comp_{\D_\kappa}(p,q)\quad\Longleftrightarrow\quad
 &\bigl[s\subseteq t\ \land\
       (\forall\xi<\dom(t))\,h(\xi)\leq t(\xi)\bigr]\notag\\
 &{}\lor\bigl[t\subseteq s\ \land\
       (\forall\xi<\dom(s))\,k(\xi)\leq s(\xi)\bigr].
 \label{eq:compat}
\end{align}
A common extension must have a stem extending both \(s\) and \(t\),
so the stems are comparable and its values on the longer stem give
the corresponding inequality. Conversely, the longer stem and the
pointwise maximum of the side functions form a common extension;
the inequality ensures normalization. All these tests use only
coordinates below the length of the longer stem. Consequently, if
both restrictions to a regular uncountable \(\beta\leq\kappa\)
are defined, then
\[
 \Comp_{\D_\kappa}(p,q)
 \quad\Longleftrightarrow\quad
 \Comp_{\D_\beta}(p\restriction\beta,q\restriction\beta).
\]
\end{remark}

\begin{lemma}\label{lem:hechler-basic}
For a regular uncountable \(\kappa\), if
\(G\subseteq\D_\kappa\) is generic, then
\(d_G\) is a function from \(\kappa\) to \(\kappa\)
that eventually dominates every function \(f\) from
\(\kappa\) to \(\kappa\) in the ground universe. Moreover,
\[
 G=\{p\in\D_\kappa:s_p\subseteq d_G
                     \text{ and }h_p\leq d_G\},
\]
where \(\D_\kappa\) on the right is the ground-model order.
In particular, \(V[G]=V[d_G]\).
\end{lemma}
\begin{proof}
Totality and eventual domination follow from the usual density
argument; see \cite[Section~4.2]{BBFM}. These facts require only
regularity of \(\kappa\). We verify the characterization of \(G\).
For \(p\in G\) and \(\xi<\kappa\), choose a common extension
in \(G\) of \(p\) and a condition whose stem includes \(\xi\).
Its value at \(\xi\) is \(d_G(\xi)\) and is at least
\(h_p(\xi)\). This proves \(s_p\subseteq d_G\) and
\(h_p\leq d_G\). Conversely, suppose a ground-model condition
\(p\) satisfies these requirements. For every \(q\in G\), the
stems of \(p\) and \(q\) are comparable, since
both are initial segments of \(d_G\). Since both side functions
are bounded by \(d_G\), the inequality in \eqref{eq:compat}
corresponding to the longer stem holds. Hence \(p\) is compatible
with every member of \(G\). The dense set of conditions extending \(p\)
or incompatible with \(p\) is met by \(G\), necessarily by a
condition extending \(p\). Therefore \(p\in G\), proving the
displayed characterization.
\end{proof}

\begin{samepage}
\begin{lemma}\label{lem:hechler-cc}
If \(\kappa\) is strongly inaccessible, then \(\D_\kappa\) is
\(\kappa^+\)-cc \cite[Section~4.2]{BBFM}.
\end{lemma}
\end{samepage}

\begin{lemma}\label{lem:hechler-closure}
For every regular uncountable \(\kappa\), the order
\(\D_\kappa\) is \(<\kappa\)-closed.
\end{lemma}
\begin{proof}
For later use, recall the lower bound witnessing closure in
\cite[Section~4.2]{BBFM}: for a decreasing sequence
\(\langle(s_i,h_i):i<\mu\rangle\), with \(0<\mu<\kappa\), put
\[
 s=\bigcup_{i<\mu}s_i,\qquad
 h(\xi)=\sup_{i<\mu}h_i(\xi)\quad(\xi<\kappa).
\]
Regularity bounds the domain and values below \(\kappa\), and
normalization fixes \(h_i(\xi)\) once \(\xi\) enters a stem.
Thus \((s,h)\) is a normalized lower bound. The empty sequence
has the greatest condition \((\varnothing,\mathbf0)\) as a
lower bound.
\end{proof}

For local domination we need predensity in \(\D_\beta\), including
conditions whose side functions are arbitrary maps from \(\beta\)
to \(\beta\). The next lemma obtains this by \(\Pi^1_1\)
reflection. In the proof of Theorem~\ref{thm:domination}, the
auxiliary ordinals \(\delta_{\alpha,i}\) will be chosen so that
\(a_{\alpha,i}\Vdash_{\D_\kappa}\delta_{\alpha,i}\in\dot E\).

\begin{lemma}\label{lem:reflection}
Let \(\kappa\) be weakly compact, let \(p\in\D_\kappa\), and
let \(\langle A_\alpha:\alpha<\kappa\rangle\) be a sequence of
maximal antichains below \(p\). Choose surjections
\(e_\alpha:\kappa\to A_\alpha\), and write
\[
 a_{\alpha,i}=e_\alpha(i)
             =(s_{\alpha,i},h_{\alpha,i})
                    \qquad(\alpha,i<\kappa).
\]
Let \(\delta:\kappa\times\kappa\to\kappa\) be any function,
and write \(\delta_{\alpha,i}=\delta(\alpha,i)\).
The set of strongly inaccessible \(\beta<\kappa\) satisfying the
following conditions is stationary in \(\kappa\):
\begin{enumerate}[label=\textup{(\roman*)}]
\item \(\ell(p)<\beta\) and \(h_p[\beta]\subseteq\beta\);
\item for all \(\alpha,i<\beta\),
      \(\ell(a_{\alpha,i})<\beta\),
      \(\delta_{\alpha,i}<\beta\), and
      \(h_{\alpha,i}[\beta]\subseteq\beta\);
\item for every \(\alpha<\beta\), the family
\[
 A_\alpha^\beta=
 \{a_{\alpha,i}\restriction\beta:i<\beta\}
\]
is predense below \(p\restriction\beta\) in \(\D_\beta\).
\end{enumerate}
\end{lemma}

\begin{proof}
Since \(\kappa\) is strongly inaccessible,
Lemma~\ref{lem:hechler-cc} gives \(1\leq|A_\alpha|\leq\kappa\).
Thus the surjections exist, with repetitions allowed. To prove
stationarity, fix a club \(B\subseteq\kappa\).
Let \(K\) be the set of uncountable fixed points of the beth
function below \(\kappa\), and put
\(K'=K\cap B\). Strong inaccessibility gives
\(\beth_\xi<\kappa\) for \(\xi<\kappa\). The beth function is
increasing and continuous, so its uncountable fixed points form a club
in \(\kappa\). Every member of \(K\) is an uncountable strong-limit
cardinal; we require regularity through a separate reflection clause.

Use finite tagged tuples to form the following predicate
\(R\subseteq V_\kappa\):
\[
\begin{aligned}
 R_0&=\{\langle0,\ell(p)\rangle\},&
 R_4&=\{\langle4,\alpha,i,\xi,h_{\alpha,i}(\xi)\rangle:
                                      \alpha,i,\xi<\kappa\},\\[.4em]
 R_1&=\{\langle1,\xi,h_p(\xi)\rangle:\xi<\kappa\},&
 R_5&=\{\langle5,\xi\rangle:\xi\in K'\},\\[.4em]
 R_2&=\{\langle2,\alpha,i,\ell(a_{\alpha,i})\rangle:
                                      \alpha,i<\kappa\},&
 R_6&=\{\langle6,s_p\rangle\},\\[.4em]
 R_3&=\{\langle3,\alpha,i,\delta_{\alpha,i}\rangle:
                                      \alpha,i<\kappa\},&
 R_7&=\{\langle7,\alpha,i,s_{\alpha,i}\rangle:
                                      \alpha,i<\kappa\},
\end{aligned}
\]
\[
 R=\bigcup_{j<8}R_j.
\]
Each stem belongs to \(V_\kappa\), since its domain and range
are bounded below the regular cardinal \(\kappa\). Side functions
need not belong to \(V_\kappa\), so we represent their values by
sections of the predicate.

We will apply reflection to the sentence
\[
 \Phi(R):=\bigwedge_{i=1}^{11}\phi_i.
\]
We describe the formulas \(\phi_1,\ldots,\phi_{11}\) below
and then verify that \(\Phi(R)\) is \(\Pi^1_1\).
In these formulas, ordinal variables range over the ordinals of the
rank structure. The symbols
\(\ell_p,\allowbreak h_p,\allowbreak\ell_{\alpha,i},\allowbreak
\delta_{\alpha,i},\allowbreak h_{\alpha,i},\allowbreak s_p,\allowbreak s_{\alpha,i}\) are read from sections \(R_0,\allowbreak R_1,\allowbreak R_2,\allowbreak R_3,
\allowbreak R_4,\allowbreak R_6,\allowbreak R_7\), respectively;
each is an abbreviation for a first-order formula using \(R\).
We use the usual first-order abbreviations for
ordinals, functions, and ordered pairs. In particular,
\(s:\eta\to\Ord\) asserts that the set \(s\) is the graph
of an ordinal-valued function with domain \(\eta\).
For a predicate \(H\), the notation \(H:\Ord\to\Ord\)
asserts that \(H\) consists of ordinal pairs and gives a
unique value at every ordinal; \(H:\eta\to\Ord\) has the
same meaning with domain exactly \(\eta\). All occurrences
of \(\Ord\) refer to the ordinals of the rank structure.
\begin{enumerate}[label=\textup{(\alph*)}]
\item The height is a limit ordinal greater than \(\omega\):
\[
\begin{aligned}
\phi_{1}\equiv (\exists w)(w=\omega)\ \land\
 (\forall\xi)(\exists\eta)\,\xi<\eta.
\end{aligned}
\]
Here \(w=\omega\) abbreviates the first-order definition
of the least infinite ordinal.
\item The sections \(R_0,\ldots,R_4\) give unique ordinal
values at every appropriate tuple of ordinal indices:
\[
\begin{aligned}
 &\phi_2 \equiv (\exists!\eta)\,\langle0,\eta\rangle\in R,\\
 &\phi_3 \equiv (\forall\xi)(\exists!\eta)\,
                         \langle1,\xi,\eta\rangle\in R,\\
 &\phi_4 \equiv (\forall\alpha)(\forall i)(\exists!\eta)\,
                         \langle2,\alpha,i,\eta\rangle\in R,\\
 &\phi_5 \equiv (\forall\alpha)(\forall i)(\exists!\eta)\,
                         \langle3,\alpha,i,\eta\rangle\in R,\\
 &\phi_6 \equiv (\forall\alpha)(\forall i)(\forall\xi)(\exists!\eta)\,
                         \langle4,\alpha,i,\xi,\eta\rangle\in R.
\end{aligned}
\]
The five lines give the values of
\(\ell_p,h_p(\xi),\ell_{\alpha,i},\delta_{\alpha,i}\),
and \(h_{\alpha,i}(\xi)\), respectively.
\item The coded stems exist as sets and are normalized:
\[
\begin{split}
 \phi_7\equiv (\exists s)\bigl[\,&\langle6,s\rangle\in R
       \ \land\ s:\ell_p\to\Ord\\
       &{}\land(\forall\xi<\ell_p)\,s(\xi)=h_p(\xi)\,\bigr],
\end{split}
\]
\[
\begin{split}
 \phi_8\equiv (\forall\alpha)(\forall i)(\exists t)\bigl[\,
       &\langle7,\alpha,i,t\rangle\in R
          \ \land\ t:\ell_{\alpha,i}\to\Ord\\
       &{}\land(\forall\xi<\ell_{\alpha,i})\,
                                 t(\xi)=h_{\alpha,i}(\xi)\\
       &{}\land s_p\subseteq t
          \ \land(\forall\xi)\,h_p(\xi)\leq h_{\alpha,i}(\xi)
       \,\bigr].
\end{split}
\]
The last line of \(\phi_8\) requires each coded condition
to extend \(p\).
\item The ordinals coded by \(R_5\) are unbounded:
\[
 \phi_9\equiv (\forall\xi)(\exists\eta)\,
       [\,\xi<\eta\ \land\ \langle5,\eta\rangle\in R\,].
\]
\item Predensity is expressed by the formula
\begin{equation*}\label{eq:predense}
\begin{split}
 \phi_{10}\equiv (\forall H)\bigl[\,(H:\Ord\to\Ord)\ \Longrightarrow
 (\forall\alpha)(\forall s)\bigl[\,
       &((\exists\eta)\,s:\eta\to\Ord)\\
       &{}\land s\subseteq H\ \land\ s_p\subseteq s\\
       &{}\land(\forall\xi)\,h_p(\xi)\leq H(\xi)\\
       &\Longrightarrow
               (\exists i)\,\Comp_R(s,H,\alpha,i)
       \,\bigr]\,\bigr].
\end{split}
\end{equation*}
The antecedent says that \((s,H)\) is normalized and extends
\(p\), with \(s\) a set in the rank structure.
The formula \(\Comp_R(s,H,\alpha,i)\) is the following
first-order test, using the stem \(t\) supplied by \(R_7\):
\[
\begin{split}
 (\exists t)\bigl[\,s_{\alpha,i}=t\ \land\
 \bigl(&[s\subseteq t\ \land\
             (\forall\xi<\dom(t))\,H(\xi)\leq t(\xi)]\\
 &{}\lor[t\subseteq s\ \land\
             (\forall\xi<\dom(s))\,
                         h_{\alpha,i}(\xi)\leq s(\xi)]\bigr)\,\bigr].
\end{split}
\]
\item Every ordinal-valued function on an ordinal has
bounded range:
\[
 \phi_{11}\equiv (\forall H)(\forall\eta)\bigl[\,
   (H:\eta\to\Ord)\ \Longrightarrow
   (\exists\zeta)(\forall\xi<\eta)\,H(\xi)<\zeta
   \,\bigr].
\]
\end{enumerate}
The formulas \(\phi_1,\ldots,\phi_9\) are first-order.
In \(\phi_{10}\) and \(\phi_{11}\), the assertions
about function graphs are first-order once \(H\) is given. Combining
the two universal second-order quantifiers gives
\begin{equation}\label{eq:reflection-sentence}
 \Phi(R)=\forall H\,\Psi(R,H),
 \qquad\Psi\text{ first-order}.
\end{equation}
Here \(R\) is the only fixed predicate; the height is not a
parameter. The side function is a universally quantified second-order
graph, while the stem is quantified as an element of the structure.

All eleven formulas hold in \((V_\kappa,\in,R)\).
The formulas \(\phi_1,\ldots,\phi_9\)
follow from the construction, \(\phi_{10}\) from maximality of \(A_\alpha\)
below \(p\) and \eqref{eq:compat}, and \(\phi_{11}\) from regularity of
\(\kappa\). By \(\Pi^1_1\)-indescribability there is a
\(\beta<\kappa\) such that
\[
 (V_\beta,\in,R\cap V_\beta)\models\Phi(R\cap V_\beta).
\]
The formula \(\phi_1\) makes \(\beta>\omega\) a limit ordinal. At a limit
height, a finite tuple belongs to \(V_\beta\) exactly when all
its entries do. Thus \(\phi_9\) implies that \(K'\cap\beta\)
is unbounded in \(\beta\). Since \(K'\) is closed,
\(\beta\in K'\subseteq K\cap B\), so \(\beta\) is an
uncountable strong-limit cardinal. For every \(\eta<\beta\),
the graph of every function from \(\eta\) to \(\beta\) is a
subset of \(V_\beta\). Under full second-order semantics,
\(\phi_{11}\) says that all these functions have bounded range.
Hence \(\operatorname{cf}(\beta)=\beta\), and \(\beta\) is
strongly inaccessible.

The reflected section \(R_0\) must contain its unique original
value \(\ell(p)\), so \(\ell(p)<\beta\). For every
\(\xi<\beta\), totality of \(R_1\) similarly requires the
tuple with value \(h_p(\xi)\), giving \(h_p(\xi)<\beta\).
The reflected sections \(R_2,R_3,R_4\) give, for all
\(\alpha,i<\beta\),
\[
 \ell(a_{\alpha,i})<\beta,\qquad
 \delta_{\alpha,i}<\beta,\qquad
 h_{\alpha,i}[\beta]\subseteq\beta.
\]
The formulas \(\phi_7\) and \(\phi_8\) supply
\(s_p\) and \(s_{\alpha,i}\) as elements
of \(V_\beta\): the stem occurs as an entry of its tuple, so a
witness from the restricted predicate is the original stem. We have
proved (i) and (ii).

For (iii), fix \(\alpha<\beta\) and a condition
\((t,H)\leq p\restriction\beta\) in \(\D_\beta\).
Regularity of \(\beta\) bounds the range of \(t\) below
\(\beta\), so \(t\in V_\beta\). The graph of
\(H:\beta\to\beta\) is a subset of \(V_\beta\). The formula \(\phi_{10}\)
therefore applies and yields \(i<\beta\) satisfying the
compatibility test with stem \(s_{\alpha,i}\). This is precisely
\eqref{eq:compat} in \(\D_\beta\), so
\(a_{\alpha,i}\restriction\beta\) is compatible with
\((t,H)\). This proves (iii), and \(\beta\in B\) proves
stationarity.
\end{proof}

\begin{definition}[Local domination]\label{def:domination}
For a regular uncountable \(\kappa\), a set
\(I\subseteq\Inacc(\kappa)\), functions
\(\vec f=\langle f_\beta:\beta\in I\rangle\) with
\(f_\beta:\beta\to\beta\), and a function \(d\) from \(\kappa\) to \(\kappa\), put
\[
 T(\vec f,d)=
 \{\beta\in I:d[\beta]\subseteq\beta
          \ \land\ f_\beta\leq^*d\restriction\beta\}.
\]
\end{definition}

\begin{theorem}\label{thm:domination}
Suppose that \(\kappa\) is weakly compact in \(V\), let
\(I=\Inacc^V(\kappa)\), and let
\(\vec f=\langle f_\beta:\beta\in I\rangle\in V\), where
\(f_\beta:\beta\to\beta\) for each \(\beta\in I\). Then
\[
 \D_\kappa\Vdash
 T(\vec f,d)\text{ is stationary in }\kappa.
\]
\end{theorem}
\begin{proof}
Work in \(V\). By Lemma~\ref{lem:hechler-closure}, names for
ordinal sequences of length less than \(\kappa\) can be decided
by a descending recursion with a lower bound. Thus \(\D_\kappa\)
adds no such sequences, preserves regularity of \(\kappa\), and
preserves the inaccessible cardinals below it.

Let \(p\in\D_\kappa\) force that \(\dot E\subseteq\kappa\)
is club. For each \(\alpha<\kappa\), every condition below
\(p\) has an extension deciding an ordinal above \(\alpha\)
that belongs to \(\dot E\), because \(p\) forces unboundedness.
Choose a maximal antichain
\(A_\alpha\) below \(p\) consisting of such conditions, enumerate
it as in Lemma~\ref{lem:reflection}, and choose
\(\delta_{\alpha,i}\) so that
\[
 \alpha<\delta_{\alpha,i}<\kappa,\qquad
 a_{\alpha,i}\Vdash\delta_{\alpha,i}\in\dot E.
\]
Choose a strongly inaccessible \(\beta\) satisfying the conclusions
of that lemma, and put \(\ell=\ell(p)\). Replace the side function
of \(p\) by
\[
 h^0(\xi)=
 \begin{cases}
 \max\{h_p(\xi),f_\beta(\xi)\},&\ell\leq\xi<\beta,\\
 h_p(\xi),&\xi<\ell\text{ or }\beta\leq\xi<\kappa.
 \end{cases}
\]
Then \(q_0=(s_p,h^0)\leq p\), and
\(h^0[\beta]\subseteq\beta\). In \(\D_\beta\), each
\(A_\alpha^\beta\), \(\alpha<\beta\), remains predense below
\(q_0\restriction\beta\).

Construct a decreasing sequence
\(\langle q_\alpha=(t_\alpha,k_\alpha):\alpha<\beta\rangle\)
with this initial condition, maintaining
\begin{equation}\label{eq:invariant}
 \dom(t_\alpha)<\beta,\qquad
 k_\alpha[\beta]\subseteq\beta,\qquad q_\alpha\leq q_0.
\end{equation}
At stage \(\alpha<\beta\), predensity gives \(i_\alpha<\beta\)
such that \(q_\alpha\restriction\beta\) is compatible with
\(a_{\alpha,i_\alpha}\restriction\beta\). By
\eqref{eq:compat}, the corresponding global conditions are
compatible as well. The longer stem and the pointwise maximum of their
side functions give a common extension with stem length below
\(\beta\) and side function mapping \(\beta\) into \(\beta\).
Extend its stem, if necessary, to
length at least \(\alpha+1\), using the same side function on
the new coordinates. Call the result \(q_{\alpha+1}\). It
satisfies \eqref{eq:invariant} and extends
\(a_{\alpha,i_\alpha}\).

At a limit \(\rho<\beta\), use the union of the previous stems
and the pointwise supremum of the previous side functions.
Regularity of \(\beta\) bounds the stem length and every value at
coordinates below \(\beta\) by \(\beta\). At coordinates at
least \(\beta\), regularity of \(\kappa\) bounds the values by
\(\kappa\). Once a coordinate belongs to a stem, its later
side-function values equal that stem value, as in
Lemma~\ref{lem:hechler-closure}; this gives normalization.
Hence the resulting condition satisfies \eqref{eq:invariant}.

Finally set
\[
 t^*=\bigcup_{\alpha<\beta}t_\alpha,\qquad
 k^*(\xi)=\sup_{\alpha<\beta}k_\alpha(\xi)
                   \quad(\xi<\kappa).
\]
The domain of \(t^*\) is exactly \(\beta\), since every stem
has length below \(\beta\) and
\(\dom(t_{\alpha+1})\geq\alpha+1\). To bound its values below
\(\beta\), fix \(\xi<\beta\) and choose \(\alpha<\beta\) with
\(\xi<\dom(t_{\alpha+1})\). The value
\(v=t_{\alpha+1}(\xi)=k_{\alpha+1}(\xi)\) is below
\(\beta\); all later side functions have this value at \(\xi\),
and all earlier ones have value at most \(v\).
Consequently
\[
 t^*(\xi)=k^*(\xi)=v<\beta.
\]
For \(\xi\geq\beta\), the supremum involves only
\(\beta<\kappa\) values, so \(k^*(\xi)<\kappa\).
Therefore \(q=(t^*,k^*)\) is a normalized condition in
\(\D_\kappa\) extending every \(q_\alpha\). It forces
\[
 d\restriction\beta=t^*,\qquad d[\beta]\subseteq\beta,
 \qquad(\forall\xi\in[\ell,\beta))\,
                       f_\beta(\xi)\leq d(\xi).
\]
Also \(q\leq a_{\alpha,i_\alpha}\) for each \(\alpha<\beta\),
so it forces all the ordinals \(\delta_{\alpha,i_\alpha}\) into
\(\dot E\). They are below \(\beta\) and cofinal in
\(\beta\), since \(\alpha<\delta_{\alpha,i_\alpha}\).
Closedness of \(\dot E\) now gives
\[
 q\Vdash\beta\in\dot E\cap T(\vec f,d).
\]
Since this construction applies below every condition forcing a
club, \(\D_\kappa\) forces \(T(\vec f,d)\) to be stationary.
\end{proof}


\begin{definition}[Limit closure points]\label{def:closure-points}
For a regular uncountable \(\kappa\) and a function \(d\) from \(\kappa\) to \(\kappa\),
put
\[
 E_d=\{\delta<\kappa:\delta>0\text{ is limit and }
                              d[\delta]\subseteq\delta\}.
\]
\end{definition}

The pressing-down lemma makes the local domination bound constant
on a stationary set of indices. A single tail \(D\) of \(E_d\)
then satisfies \(D\cap\beta\subseteq C_\beta\) at these indices.

\begin{corollary}\label{cor:capture}
Suppose that \(\kappa\) is weakly compact in \(V\), let
\(I=\Inacc^V(\kappa)\), and let
\(\langle C_\beta:\beta\in I\rangle\in V\) consist of clubs
\(C_\beta\subseteq\beta\). If \(G\subseteq\D_\kappa\) is
generic over \(V\) and \(d=d_G\), then in \(V[d]\) there are
a club \(D\subseteq\kappa\) and a stationary set
\(T\subseteq I\) such that
\[
 D\cap\beta\subseteq C_\beta\qquad(\beta\in T).
\]
In particular, every ground-model \(C\)-sequence whose domain
contains \(I\) is non-amenable in \(V[d]\).
\end{corollary}
\begin{proof}
In \(V\), for each \(\beta\in I\) form the next-point function
\[
 n_\beta(\xi)=\min(C_\beta\setminus(\xi+1))
                              \qquad(\xi<\beta).
\]
Unboundedness gives \(\xi<n_\beta(\xi)<\beta\). Apply
Theorem~\ref{thm:domination} to these functions. In \(V[d]\),
\(\kappa\) is regular and
\[
 S=\{\beta\in I:d[\beta]\subseteq\beta
                  \ \land\ n_\beta\leq^*d\restriction\beta\}
\]
is stationary. For \(\beta\in S\), let
\(\gamma_\beta<\beta\) be the least bound witnessing
\(n_\beta\leq^*d\restriction\beta\). The pressing-down lemma, applied
in \(V[d]\), gives a stationary \(T\subseteq S\) and an
ordinal \(\gamma<\kappa\) such that
\(\gamma_\beta=\gamma\) for all \(\beta\in T\).

The set \(E_d\) is closed: if \(\delta<\kappa\) is a limit
of its members, then for each \(\xi<\delta\) there is
\(\varepsilon\in E_d\cap\delta\) with \(\xi<\varepsilon\),
giving \(d(\xi)<\varepsilon<\delta\). It is also unbounded.
Above any prescribed bound, choose a strictly increasing sequence
\(\langle\eta_n:n<\omega\rangle\) such that
\(\sup(d[\eta_n])<\eta_{n+1}\). Regularity of \(\kappa\)
permits each choice and gives
\(\delta=\sup_{n<\omega}\eta_n<\kappa\); this \(\delta\)
belongs to \(E_d\). Therefore
\(D=E_d\setminus(\gamma+1)\) is club.

Fix \(\beta\in T\) and \(\delta\in D\cap\beta\). Then
\[
 \xi<n_\beta(\xi)\leq d(\xi)<\delta
                     \qquad(\gamma\leq\xi<\delta).
\]
For every \(\eta<\delta\), choose
\(\xi\in[\max\{\eta,\gamma\},\delta)\); this is possible
because \(\delta>\gamma\) is limit. The point
\(n_\beta(\xi)\) lies in \(C_\beta\cap\delta\) and is
greater than \(\eta\). Hence \(C_\beta\cap\delta\) is
unbounded in \(\delta\), and closedness of \(C_\beta\) gives
\(\delta\in C_\beta\). This proves the asserted inclusion.
For a \(C\)-sequence with domain containing \(I\), the resulting
capture set contains \(T\) and is therefore stationary in
\(V[d]\).
\end{proof}

\section{Hechler iterations}
\label{sec:iteration}

We present \(<\kappa\)-support iterations as partial functions
whose domains are their supports. Stronger conditions are smaller
in the order. Unless specified otherwise, the iteration and all
cardinality calculations are computed in a fixed ground universe
\(V\). At each stage the iterand is the full normalized Hechler
order of the intermediate extension. The decided-stem,
\(\Delta\)-system, and name-counting arguments below adapt
\cite[Lemma 3.10]{Schuerz}, omitting the second forcing factor
used there.

\begin{definition}[The Hechler iteration]\label{def:iteration}
Let \(\kappa\) be strongly inaccessible and let \(\lambda\) be
an ordinal. Define the \(<\kappa\)-support iteration
\[
 \langle\PP_\alpha:\alpha\leq\lambda\rangle,
 \qquad\langle\dot{\mathbb Q}_\alpha:\alpha<\lambda\rangle
\]
recursively, with \(\PP_0=\{\varnothing\}\) and
\[
 \mathbf1_\alpha\Vdash_{\PP_\alpha}
 \dot{\mathbb Q}_\alpha=\D_\kappa^{V[G_\alpha]}.
\]
The greatest condition of each \(\PP_\alpha\) is the empty
function, denoted by \(\mathbf1_\alpha\). The name for the
greatest condition of the iterand is
\(\dot{\mathbf1}_\alpha=\check{(\varnothing,\mathbf0)}\),
where \(\mathbf0\) is the constant zero function from
\(\kappa\) to \(\kappa\).

Choose \(\dot{\mathbb Q}_\alpha\) to be a \emph{full} name:
for every \(r\in\PP_\alpha\) and every
\(\PP_\alpha\)-name \(\tau\),
\begin{equation}\label{eq:full-iterand-name}
 r\Vdash\tau\in\dot{\mathbb Q}_\alpha
 \quad\Longrightarrow\quad
 (\exists\sigma\in\dom(\dot{\mathbb Q}_\alpha))\,
                         r\Vdash\sigma=\tau.
\end{equation}
Here \(\dom(\dot{\mathbb Q}_\alpha)\) is the set of names
occurring as first entries in \(\dot{\mathbb Q}_\alpha\).
Such a name can be chosen with a set-sized domain. To see this,
mix a name forced by \(r\) to be a Hechler condition with
\(\dot{\mathbf1}_\alpha\) on the other members of a maximal
antichain containing \(r\). The mixed name is globally forced
to be a Hechler condition. Its two component graphs have names
whose entries are check names for pairs in \(\kappa\times\kappa\),
paired with conditions of \(\PP_\alpha\). These graph names form
a set. Take all ordered-pair names formed from them that are
globally forced to be normalized Hechler conditions, each paired
with \(\mathbf1_\alpha\), and include
\(\dot{\mathbf1}_\alpha\). Every Hechler condition in an extension
has such a representative, and the mixed representative agrees
with the original name below \(r\). This gives
\eqref{eq:full-iterand-name} without further strengthening \(r\).

For \(0<\theta\leq\lambda\), a condition \(p\in\PP_\theta\)
is a partial function with
\[
 \supp(p)=\dom(p)\subseteq\theta,
 \qquad |\supp(p)|<\kappa,
\]
such that \(p\restriction\alpha\in\PP_\alpha\) for every
\(\alpha<\theta\), and, at each \(\alpha\in\dom(p)\),
\[
 p(\alpha)\in\dom(\dot{\mathbb Q}_\alpha),
 \qquad p\restriction\alpha\Vdash_{\PP_\alpha}
                         p(\alpha)\in\dot{\mathbb Q}_\alpha.
\]
Membership in the iterand is required below
\(p\restriction\alpha\).
Whenever \(\alpha\notin\dom(p)\), the expression
\(p(\alpha)\) is understood to mean
\(\dot{\mathbf1}_\alpha\); this convention does not add
\(\alpha\) to the domain. Initial segments are included by
leaving all later coordinates omitted. At a limit stage we
take all such partial functions whose restrictions belong
to the earlier stages, including unbounded supports when
the cofinality of the stage is less than \(\kappa\).

For \(p,q\in\PP_\theta\), the order is
\begin{equation}\label{eq:iteration-order}
 \begin{split}
 q\leq p\quad\Longleftrightarrow\quad&
 \supp(p)\subseteq\supp(q)\ \text{ and}\\
 & (\forall\alpha\in\supp(p))\,
 q\restriction\alpha\Vdash_{\PP_\alpha}
 q(\alpha)\leq_{\dot{\mathbb Q}_\alpha}p(\alpha).
 \end{split}
\end{equation}
At a coordinate added by \(q\), validity of \(q\) already
ensures that its value extends the greatest condition assigned
to that coordinate in \(p\).
A coordinate in \(\dom(p)\) remains in the support even when
its value is forced to be the greatest condition. Thus one may
insert coordinates with trivial values, and strengthening
retains every coordinate in the domain.
These presentations are preorders; identifying mutually extending
conditions gives partial orders.
\end{definition}

\begin{definition}[Decided stems]\label{def:decided-stems}
For \(p\in\PP_\lambda\) and \(\alpha<\lambda\), choose
\(\PP_\alpha\)-names \(\dot s_\alpha^p\) and
\(\dot h_\alpha^p\) with
\[
 p\restriction\alpha\Vdash
 p(\alpha)=\langle\dot s_\alpha^p,\dot h_\alpha^p\rangle.
\]
Below this initial condition, the names denote a stem in
\({}^{<\kappa}\kappa\) and a function in \({}^{\kappa}\kappa\),
respectively. At omitted coordinates use the empty stem and the
zero function. Any two choices of component names agree below
\(p\restriction\alpha\), so the following definition is
independent of those choices.
Define
\[
 \mathcal E_\lambda=
 \{p\in\PP_\lambda:(\forall\alpha\in\supp(p))
   (\exists s_\alpha\in{}^{<\kappa}\kappa)\,
   p\restriction\alpha\Vdash
                   \dot s_\alpha^p=\check s_\alpha\}.
\]
Members of \(\mathcal E_\lambda\) have \emph{decided stems};
their side-function names remain arbitrary.
\end{definition}

\begin{lemma}\label{lem:factorization}
Let \(\PP_\lambda\) be the Hechler iteration at a strongly
inaccessible \(\kappa\), and let \(\theta\leq\lambda\).
The inclusion \(\PP_\theta\subseteq\PP_\lambda\) preserves
incompatibility and sends maximal antichains to predense sets.
Consequently, if \(G_\lambda\) is generic over \(V\), then
\(G_\theta=G_\lambda\cap\PP_\theta\) is
\(\PP_\theta\)-generic over \(V\).
If \(\PP_\theta\) preserves \(\kappa\) and adds no ordinal
sequences of length less than \(\kappa\), then, in
\(V[G_\theta]\), the remaining forcing is forcing equivalent
to the full \(<\kappa\)-support Hechler iteration on
\([\theta,\lambda)\), computed in that universe.
\end{lemma}
\begin{proof}
The case \(\theta=\lambda\) is immediate. Suppose
\(\theta<\lambda\). Initial-segment completeness and genericity
are the standard iteration facts recalled in
\cite[Section~7, pp.~799--800]{Cummings}. In our partial-function
presentation, the restriction map is a projection: if
\(p\in\PP_\lambda\) and \(r\leq p\restriction\theta\), then
\begin{equation}\label{eq:iteration-splice}
 r\mathbin{\smallfrown}_\theta p
       =r\cup(p\restriction[\theta,\lambda))
\end{equation}
is a condition below both \(p\) and \(r\).

By \cite[Proposition~7.11]{Cummings}, the quotient is forcing
equivalent to the factor iteration obtained by translating the
remaining iterand names into \(V[G_\theta]\). Fullness,
\eqref{eq:full-iterand-name}, permits choosing coordinate
representatives in our iterand-name domains without further
strengthening. At successor stages the translated iterands are
the full Hechler orders of the corresponding intermediate
extensions.

For the limit stages, the hypothesis gives
\[
 ([\lambda]^{<\kappa})^{V[G_\theta]}
                         =([\lambda]^{<\kappa})^V,
\]
since an enumeration of a set on the left is an ordinal
sequence of length less than \(\kappa\). More generally,
there are no new sequences of fewer than \(\kappa\)
ground-model objects. Thus the limit-stage construction of the
factor iteration admits exactly the \(<\kappa\)-supports of
the full tail, including unbounded supports at stages of
cofinality less than \(\kappa\). Consequently,
\[
 \PP_\lambda\simeq
           \PP_\theta*\dot\PP_{[\theta,\lambda)},
\]
where the second factor is the full \(<\kappa\)-support Hechler
iteration computed in \(V[G_\theta]\).
\end{proof}

\begin{lemma}\label{lem:iteration}
Let \(\kappa\) be strongly inaccessible. For every ordinal
\(\lambda\), the Hechler iteration \(\PP_\lambda\) is
\(<\kappa\)-closed and \(\kappa^+\)-Knaster, and
\(\mathcal E_\lambda\) is dense in \(\PP_\lambda\).
Moreover, for every \(\theta\leq\lambda\),
\[
 \PP_\theta\Vdash
 \text{the tail iteration }\PP_{[\theta,\lambda)}
 \text{ is }{<}\kappa\text{-closed},
\]
where the tail is computed in the initial generic extension
as in Lemma~\ref{lem:factorization}.
\end{lemma}

\begin{proof}
\emph{Closure and tails.}
We prove closure and preservation of strong inaccessibility
simultaneously by induction on \(\lambda\). Let
\(\langle p_i:i<\delta\rangle\) be a decreasing sequence in
\(\PP_\lambda\), where \(0<\delta<\kappa\), and put
\[
 u=\bigcup_{i<\delta}\supp(p_i),\qquad |u|<\kappa.
\]
Construct \(q\) recursively on \(u\). Suppose that
\(r_\alpha=q\restriction\alpha\leq
p_i\restriction\alpha\) for every \(i<\delta\).
Choose \(\PP_\alpha\)-names \(\dot t_\alpha\) and
\(\dot g_\alpha\) such that
\[
 r_\alpha\Vdash
 \dot t_\alpha=\bigcup_{i<\delta}\dot s_\alpha^{p_i},
 \qquad
 r_\alpha\Vdash
 \dot g_\alpha=
       \left\langle\sup_{i<\delta}\dot h_\alpha^{p_i}(\xi):
                                     \xi<\kappa\right\rangle.
\]
Choose a name \(\dot b_\alpha\) with
\[
 r_\alpha\Vdash
 \dot b_\alpha=\langle\dot t_\alpha,\dot g_\alpha\rangle.
\]
By the induction hypothesis \(\kappa\) remains regular
after \(\PP_\alpha\). Under \(r_\alpha\), the coordinate values
form a decreasing sequence of Hechler conditions, with the maximum
used at omitted coordinates. Lemma~\ref{lem:hechler-closure} gives
\[
 r_\alpha\Vdash
 \dot b_\alpha\in\dot{\mathbb Q}_\alpha
       \ \land\ (\forall i<\delta)\,
                           \dot b_\alpha\leq p_i(\alpha).
\]
By \eqref{eq:full-iterand-name}, choose
\(q(\alpha)\in\dom(\dot{\mathbb Q}_\alpha)\) with
\(r_\alpha\Vdash q(\alpha)=\dot b_\alpha\).
Set \(\dom(q)=u\), omitting all coordinates outside \(u\).
Induction using \eqref{eq:iteration-order} yields
\(q\in\PP_\lambda\) and \(q\leq p_i\) for all
\(i<\delta\). Denote this lower bound by
\(\operatorname{lb}_\lambda(\langle p_i:i<\delta\rangle)\).
Its support is exactly \(u\), including coordinates whose
lower-bound value is trivial. At each coordinate its initial
restriction forces its stem to be the union of the corresponding
stems. For an empty sequence take
\(\mathbf1_\lambda\).

For a name \(\dot f:\eta\to\Ord\), with \(\eta<\kappa\),
and a condition forcing its stated domain, construct
\(\langle r_\xi:\xi\leq\eta\rangle\) below that condition
with
\[
 r_{\xi+1}\Vdash\dot f(\xi)=\check\beta_\xi,
 \qquad
 r_\zeta=\operatorname{lb}_\lambda
                   (\langle r_\xi:\xi<\zeta\rangle)
                  \quad(0<\zeta\leq\eta\text{ limit}).
\]
Then
\(r_\eta\Vdash\dot f=
\check{\langle\beta_\xi:\xi<\eta\rangle}\).
Thus no ordinal sequences of length less than \(\kappa\)
are added. In particular, no subsets of ordinals below
\(\kappa\) are added, no cardinal below \(\kappa\) is
collapsed, and \(\kappa\) remains regular. It therefore remains
strongly inaccessible, completing the induction.
Lemma~\ref{lem:factorization}
identifies the quotient after \(\PP_\theta\) with the full
Hechler tail in \(V[G_\theta]\); the same lower-bound
recursion in that universe proves its \(<\kappa\)-closure.

\emph{Density of decided stems.}
Fix \(p\in\PP_\lambda\). Starting with \(p_0=p\), construct
a decreasing sequence \(\langle p_n:n<\omega\rangle\).
At stage \(n\), enumerate
\[
 \supp(p_n)=\{\alpha_{n,\xi}:\xi<\mu_n\},
                         \qquad\mu_n<\kappa,
\]
and build decreasing conditions
\(\langle p_{n,\xi}:\xi\leq\mu_n\rangle\), starting with
\(p_{n,0}=p_n\). At \(\xi<\mu_n\), let
\(\alpha=\alpha_{n,\xi}\). First strengthen
\(p_{n,\xi}\restriction\alpha\) to decide the stem length.
Then use \(<\kappa\)-closure of \(\PP_\alpha\) to decide
all values of that stem. This gives
\(v_{n,\xi}\leq p_{n,\xi}\restriction\alpha\) and
\(t_{\alpha,n}\in{}^{<\kappa}\kappa\) such that
\[
 v_{n,\xi}\Vdash
                    \dot s_\alpha^{p_{n,\xi}}=\check t_{\alpha,n}.
\]
Set
\[
 p_{n,\xi+1}=v_{n,\xi}
                   \cup(p_{n,\xi}\restriction[\alpha,\lambda)).
\]
By \eqref{eq:iteration-splice}, this is a condition below
\(p_{n,\xi}\), and its initial restriction forces its stem
at \(\alpha\) to be \(t_{\alpha,n}\). At a nonzero limit
\(\zeta\leq\mu_n\), use the specified lower bound
\[
 p_{n,\zeta}=\operatorname{lb}_\lambda
                   (\langle p_{n,\xi}:\xi<\zeta\rangle).
\]
Finally put \(p_{n+1}=p_{n,\mu_n}\).
Supports increase throughout the construction, and the lower-bound
operation takes their union. Hence every coordinate introduced
during stage \(n\) belongs to \(\supp(p_{n+1})\) and is
treated at every later stage. Each stage involves fewer than
\(\kappa\) steps, so all its conditions have support of size less
than \(\kappa\).

Take
\[
 q=\operatorname{lb}_\lambda(\langle p_n:n<\omega\rangle),
 \qquad \supp(q)=\bigcup_{n<\omega}\supp(p_n).
\]
Fix \(\alpha\in\supp(q)\), and choose \(n_\alpha\) so
that \(\alpha\in\supp(p_n)\) for all \(n\geq n_\alpha\).
The stem decision at stage \(n\), together with the decreasing order,
gives
\[
 q\restriction\alpha\Vdash
 \dot s_\alpha^{p_n}\subseteq\check t_{\alpha,n}
                     \subseteq\dot s_\alpha^{p_{n+1}}
                               \quad(n\geq n_\alpha).
\]
In particular, \(t_{\alpha,n}\subseteq t_{\alpha,n+1}\).
Regularity of \(\kappa\) gives
\(t_\alpha=\bigcup_{n\geq n_\alpha}t_{\alpha,n}
\in{}^{<\kappa}\kappa\). The two inclusions above show that
the union of the stems of the \(p_n\)'s is exactly
\(t_\alpha\); stems from earlier stages are contained in
\(t_{\alpha,n_\alpha}\). By the choice of lower bound,
\[
 q\restriction\alpha\Vdash
 \dot s_\alpha^q=\bigcup_{n<\omega}\dot s_\alpha^{p_n}
                                      =\check t_\alpha.
\]
Thus \(q\in\mathcal E_\lambda\) and \(q\leq p\).

\emph{The Knaster property.}
Given \(\langle p_i:i<\kappa^+\rangle\), choose
\(q_i\in\mathcal E_\lambda\) with \(q_i\leq p_i\),
and put \(u_i=\supp(q_i)\). Since
\(\kappa^{<\kappa}=\kappa\), the generalized
\(\Delta\)-system lemma gives \(J\subseteq\kappa^+\)
and a root \(R\) with
\[
 |J|=\kappa^+,\qquad |R|<\kappa,\qquad
 u_i\cap u_j=R\quad(i\neq j\text{ in }J).
\]
If one support occurs \(\kappa^+\) times, use it as the root;
otherwise there are \(\kappa^+\) distinct supports to which the
lemma applies. The number of assignments
of decided stems on \(R\) is
\[
 \left|{}^R({}^{<\kappa}\kappa)\right|
                              =\kappa^{|R|}\leq\kappa.
\]
Thus there is \(J'\subseteq J\) of size \(\kappa^+\)
and stems \(s_\alpha\), \(\alpha\in R\), with
\[
 q_i\restriction\alpha\Vdash
                      \dot s_\alpha^{q_i}=\check s_\alpha
                              \quad(i\in J',\ \alpha\in R).
\]

For distinct \(i,j\in J'\), construct \(w\) recursively
on \(u_i\cup u_j\), retaining all these coordinates.
Off the root use the name from the unique condition whose
support contains the coordinate; omit coordinates outside
\(u_i\cup u_j\). Inductively, suppose
\(w\restriction\alpha\leq
q_i\restriction\alpha,q_j\restriction\alpha\).
At a root coordinate choose names \(\dot g_\alpha\) and
\(\dot c_\alpha\) such that
\[
 w\restriction\alpha\Vdash
 \dot g_\alpha=\left\langle
       \max\{\dot h_\alpha^{q_i}(\xi),\dot h_\alpha^{q_j}(\xi)\}:
                                      \xi<\kappa\right\rangle,
\]
\[
 w\restriction\alpha\Vdash
 \dot c_\alpha=\langle\check s_\alpha,\dot g_\alpha\rangle.
\]
This initial restriction forces both stems to be \(s_\alpha\).
Normalization implies that both side functions agree with
\(s_\alpha\) on its domain. Their pointwise maximum therefore
does as well, and
\[
 w\restriction\alpha\Vdash
 \dot c_\alpha\in\dot{\mathbb Q}_\alpha
              \ \land\ \dot c_\alpha\leq q_i(\alpha),q_j(\alpha).
\]
Fullness gives \(w(\alpha)\in\dom(\dot{\mathbb Q}_\alpha)\)
with \(w\restriction\alpha\Vdash
w(\alpha)=\dot c_\alpha\). Outside the root, the retained name
is valid below the stronger initial restriction and extends the
greatest condition assigned to that coordinate in the other
condition. Since
\(|u_i\cup u_j|<\kappa\), the recursion defines
\(w\in\PP_\lambda\) with \(w\leq q_i,q_j\),
and hence \(w\leq p_i,p_j\). This proves
\(\kappa^+\)-Knaster.

To record the cardinal preservation needed below, let
\(\dot f:\eta\to\kappa^+\) be a name, with \(\eta\leq\kappa\).
Choose maximal antichains deciding its values and let
\(X_\xi\) be the set of possible values at \(\xi<\eta\).
Then
\[
 |X_\xi|\leq\kappa,\qquad
 \left|\bigcup_{\xi<\eta}X_\xi\right|\leq\kappa,
 \qquad
 \Vdash\operatorname{ran}(\dot f)
                           \subseteq\bigcup_{\xi<\eta}X_\xi.
\]
The union is bounded in the ground-model \(\kappa^+\), so
no such map is cofinal. Since \(\kappa\) is preserved, the
ground-model \(\kappa^+\) remains its regular successor.
\end{proof}

\begin{lemma}\label{lem:iteration-size}
Let \(\kappa\) be strongly inaccessible and put \(\mu=2^\kappa\).
For every \(\theta\leq\kappa^+\), the Hechler iteration
\(\PP_\theta\) has a \(<\kappa\)-closed dense suborder of
size at most \(\mu\).
\end{lemma}
\begin{proof}
All counts are in \(V\). Cardinal arithmetic gives
\[
 \mu^\kappa=\mu,\qquad \mu^{<\kappa}=\mu,
 \qquad\kappa^+\leq\mu.
\]
Construct dense sets \(D_\theta\subseteq\PP_\theta\)
recursively, starting with \(D_0=\PP_0\). Suppose
\(D_\alpha\) has been constructed for each \(\alpha<\theta\)
and \(|D_\alpha|\leq\mu\).

For \(\alpha<\theta\), let \(H_\alpha\) be the set of
names for graphs of functions of the form
\[
 \dot h_{\vec A,\vec\nu}=
 \{(\check{\langle\xi,\nu_\xi(a)\rangle},a):
                       \xi<\kappa,\ a\in A_\xi\},
\]
where each \(A_\xi\subseteq D_\alpha\) is a maximal
antichain of \(\PP_\alpha\) and
\(\nu_\xi:A_\xi\to\kappa\). Such a name is globally
forced to be a function from \(\kappa\) to \(\kappa\).
Conversely, for every globally total function name
\(\dot h\), density of \(D_\alpha\) gives maximal
antichains \(A_\xi\subseteq D_\alpha\) and assignments with
\[
 a\Vdash\dot h(\xi)=\check{\nu_\xi(a)}
                            \quad(\xi<\kappa,\ a\in A_\xi).
\]
The corresponding member of \(H_\alpha\) is globally
forced equal to \(\dot h\). By the chain condition,
each \(A_\xi\) has size at most \(\kappa\). Thus each
presentation is coded by at most \(\kappa\) triples in
\(\kappa\times D_\alpha\times\kappa\), and
\[
 |H_\alpha|\leq
            (\kappa\cdot\mu\cdot\kappa)^\kappa=\mu.
\]

For \(s\in{}^{<\kappa}\kappa\) and \(\dot h\in H_\alpha\),
define the normalized function name \(\dot h[s]\) by
\[
 \mathbf1_\alpha\Vdash\dot h[s](\xi)=
 \begin{cases}
  s(\xi),&\xi<\dom(s),\\
  \dot h(\xi),&\dom(s)\leq\xi<\kappa.
 \end{cases}
\]
Then \(\mathbf1_\alpha\Vdash
\langle\check s,\dot h[s]\rangle
                                      \in\dot{\mathbb Q}_\alpha\).
For each such \(s,\dot h\), fullness gives a name
\(\sigma\in\dom(\dot{\mathbb Q}_\alpha)\); fix one with
\[
 \mathbf1_\alpha\Vdash
                 \sigma=\langle\check s,\dot h[s]\rangle.
\]
Let \(F_\alpha\) consist of these chosen names together
with \(\dot{\mathbf1}_\alpha\).
Since \(|{}^{<\kappa}\kappa|=\kappa\),
\(|F_\alpha|\leq\mu\). Define
\[
 D_\theta=\{p\in\PP_\theta:
       (\forall\alpha\in\supp(p))\ p(\alpha)\in F_\alpha\}.
\]

To prove density, fix \(p\in\PP_\theta\) and choose
\(p'\in\mathcal E_\theta\) with \(p'\leq p\).
For \(\alpha\in\supp(p')\), let \(s_\alpha\) be its
decided stem. The component name \(\dot h_\alpha^{p'}\)
is forced to be a function only below
\(p'\restriction\alpha\). First extend it to a globally
total function-name by mixing with the zero function off
that cone: use a maximal antichain containing
\(p'\restriction\alpha\), retain the original name there,
and use zero on the other members. Then choose its nice
representative \(\dot h_\alpha\in H_\alpha\). Thus
\[
 p'\restriction\alpha\Vdash
                         \dot h_\alpha=\dot h_\alpha^{p'}.
\]
Normalization of \(p'(\alpha)\) gives
\[
 p'\restriction\alpha\Vdash
 p'(\alpha)=\langle\check s_\alpha,\dot h_\alpha[s_\alpha]\rangle.
\]
At every \(\alpha\in\supp(p')\), use for \(q(\alpha)\)
the name in \(F_\alpha\) chosen for
\(s_\alpha,\dot h_\alpha\), retaining that coordinate
even if its value is trivial. Omit coordinates outside
\(\supp(p')\). Thus \(\supp(q)=\supp(p')\).
Induction on the coordinates gives
\[
 q\restriction\alpha\leq p'\restriction\alpha,
 \qquad q\restriction\alpha\Vdash q(\alpha)=p'(\alpha).
\]
Hence \(q\in D_\theta\) and \(q\leq p'\leq p\).
Finally,
\[
 |[\theta]^{<\kappa}|\leq\mu^{<\kappa}=\mu,
 \qquad |D_\theta|\leq\mu\cdot\mu^{<\kappa}=\mu.
\]
This completes the recursion at successor and limit lengths.
For a descending sequence of length less than \(\kappa\) in
\(D_\theta\), Lemma~\ref{lem:iteration} gives a lower bound
in \(\PP_\theta\). Refine that lower bound into \(D_\theta\)
by density. Thus \(D_\theta\), with the inherited order, is
\(<\kappa\)-closed.
\end{proof}

\begin{lemma}\label{lem:localization}
Let \(\kappa\) be strongly inaccessible, let \(\lambda=\kappa^+\),
and let \(G_\lambda\subseteq\PP_\lambda\) be generic over
\(V\) for the Hechler iteration. Put
\(G_\theta=G_\lambda\cap\PP_\theta\) for
\(\theta\leq\lambda\). Every subset of \(\kappa\) in
\(V[G_\lambda]\) belongs to \(V[G_\theta]\) for some
\(\theta<\lambda\). The same holds for every \(C\)-sequence
on \(\kappa\).
\end{lemma}
\begin{proof}
Let \(X\subseteq\kappa\) belong to \(V[G_\lambda]\).
Choose a name \(\dot X\) that is globally forced to
be a subset of \(\kappa\) and evaluates to \(X\).
Such a name is obtained by intersecting any name for
\(X\) with \(\check\kappa\). For each \(\xi<\kappa\),
choose a maximal antichain \(B_\xi\) deciding whether
\(\check\xi\in\dot X\), and set
\[
 B_\xi^+=\{p\in B_\xi:p\Vdash\check\xi\in\dot X\},
 \qquad
 \dot X^*=\{(\check\xi,p):\xi<\kappa,\ p\in B_\xi^+\}.
\]
Then \(\mathbf1_\lambda\Vdash\dot X^*=\dot X\).
By Lemma~\ref{lem:iteration}, \(|B_\xi|\leq\kappa\).
The union
\[
 a=\bigcup_{\xi<\kappa}\ \bigcup_{p\in B_\xi}\supp(p)
\]
has size at most \(\kappa\). Regularity of
\(\lambda=\kappa^+\) gives
\[
 \theta=\sup\{\alpha+1:\alpha\in a\}<\lambda.
\]
Every condition used in \(\dot X^*\) has support
below \(\theta\) and is a condition of \(\PP_\theta\):
at each coordinate \(\alpha\), its entry is already
a \(\PP_\alpha\)-name. Hence \(\dot X^*\) is a
\(\PP_\theta\)-name and
\[
 X=(\dot X^*)^{G_\lambda}
        =(\dot X^*)^{G_\theta}\in V[G_\theta].
\]

Fix in \(V\) a bijection \(b:\kappa\times\kappa\to\kappa\).
For a \(C\)-sequence \(\vec C\) in the final
extension, apply the first part to
\[
 X_{\vec C}=\{b(\beta,\xi):\beta<\kappa,\ \xi\in C_\beta\}.
\]
In the resulting intermediate extension, define
\[
 C_\beta=\{\xi<\beta:b(\beta,\xi)\in X_{\vec C}\}
                                             \qquad(\beta<\kappa).
\]
These are the original entries. For a fixed set of
ordinals, closedness and unboundedness in a fixed ordinal
are absolute between the two extensions. Thus the entire
sequence is already a \(C\)-sequence in \(V[G_\theta]\).
\end{proof}

\begin{lemma}\label{lem:stationarity}
Let \(\kappa\) be regular uncountable, let \(S\subseteq\kappa\)
be stationary, and let \(\mathbb R\) be \(<\kappa\)-closed.
Then \(\mathbb R\) preserves the stationarity of \(S\).
\end{lemma}
\begin{proof}
Closure adds no ordinal sequences of length \(<\kappa\)
and hence preserves regularity of \(\kappa\). Fix
\(p\in\mathbb R\) with
\(p\Vdash\dot E\subseteq\kappa\text{ is club}\).
In the ground universe construct decreasing conditions
\(\langle p_i:i<\kappa\rangle\), a strictly increasing
continuous sequence \(\langle\eta_i:i<\kappa\rangle\),
and ordinals \(\xi_i<\kappa\), starting with
\(p_0=p\) and \(\eta_0=0\). At successors, forced
unboundedness lets us choose
\[
 p_{i+1}\leq p_i,\qquad
 \eta_i<\xi_i<\kappa,\qquad
 p_{i+1}\Vdash\check\xi_i\in\dot E,
 \qquad\eta_{i+1}=\xi_i+1.
\]
At a nonzero limit \(\rho<\kappa\), choose a lower
bound \(p_\rho\) of the preceding conditions and put
\(\eta_\rho=\sup_{i<\rho}\eta_i<\kappa\).
Continuity and strict increase make
\[
 F=\{\eta_\rho:0<\rho<\kappa\text{ is a limit ordinal}\}
\]
a club in \(\kappa\). Choose a limit \(\rho\) with
\(\eta_\rho\in S\). Since \(p_\rho\) extends
every \(p_{i+1}\) for \(i<\rho\),
\[
 p_\rho\Vdash\check{\{\xi_i:i<\rho\}}\subseteq\dot E,
 \qquad
 \sup_{i<\rho}\xi_i=\sup_{i<\rho}\eta_{i+1}=\eta_\rho.
\]
The forced closedness of \(\dot E\) gives
\(p_\rho\Vdash\check\eta_\rho\in\dot E\cap\check S\).
Repeating the construction below any extension of
\(p\) shows that conditions forcing an intersection
with \(S\) are dense below \(p\), as required.
\end{proof}


\section{Preservation and the consistency result}\label{sec:proof-main}
\label{sec:elimination}

\subsection{Preservation under the \texorpdfstring{\(\Pi^2_1\)}{Pi-2-1} hypothesis}

For the \(\Pi^2_1\) construction, choose the local Menas function
\(f\) at height \(\kappa+1\) specified in
Proposition~\ref{prop:target}. Let \(\RR\) be the local lottery preparation of
\cite{Johnstone}, relative to \(f\). This is the Easton support iteration
of length \(\kappa\) which, at eligible stages
\(\gamma<\kappa\), chooses from the lottery sum of
\(<\gamma\)-closed orders in
\(H_{f(\gamma)^+}\), calculated in the intermediate
universe. Eligibility requires
\(\gamma\in\dom(f)\) and \(f''\gamma\subseteq\gamma\).
We use the size-\(\kappa\) presentation with conditions
in \(V_\kappa\) described in
\cite[pp.~1229--1230]{Johnstone}. It is \(\kappa\)-cc and
preserves inaccessibility. By
\cite[Corollary 38]{Johnstone}, with \(m=1\), it also
preserves \(\Pi^2_1\)-indescribability. The
indestructibility bound at this height is
\(\beth_\kappa=\kappa\).

\begin{theorem}
\label{thm:preparation}
Suppose \(\kappa\) is \(\Pi^2_1\)-indescribable and
\(2^\kappa=\kappa^+\). Let \(G\subseteq\RR\) be
generic for the local lottery preparation above and put
\(U=V[G]\). Then \(U\models2^\kappa=\kappa^+\), and,
for every \(\theta\leq\kappa^+\), the full
\(<\kappa\)-support Hechler iteration computed in \(U\)
satisfies
\[
 \PP_\theta\Vdash_U
          \text{\(\kappa\) is \(\Pi^2_1\)-indescribable}.
\]
\end{theorem}

The proof is deferred to Appendix~\ref{sec:preservation}.
It uses the canonical presentation of
Subsection~\ref{subsec:canonical} to extend
\cite[Corollary 38]{Johnstone}, with \(m=1\), to the full
Hechler iteration.

\begin{corollary}
\label{cor:iteration-preservation}
In the universe prepared in
Theorem~\ref{thm:preparation}, all initial segments
and the full Hechler iteration of length
\(\kappa^+\) preserve
\(\Pi^2_1\)-indescribability and, in particular,
weak compactness.
\end{corollary}
\begin{proof}
Apply Theorem~\ref{thm:preparation} to each
initial segment computed in the prepared universe.
\end{proof}

\subsection{Preservation from a \texorpdfstring{\(\Pi^3_1\)}{Pi-3-1} hypothesis}
\label{subsec:pi3-route}

The stronger hypothesis of \(\Pi^3_1\)-indescribability allows
the iteration to be treated directly by
\cite[Corollary 38]{Johnstone}. This application requires the local lottery preparation
at height \(\kappa+2\), rather than the preparation at height
\(\kappa+1\) used in Theorem~\ref{thm:preparation}.

\begin{proposition}
\label{prop:pi3-preparation}
Suppose that \(\kappa\) is \(\Pi^3_1\)-indescribable and
\(2^\kappa=\kappa^+\). Choose a partial local Menas function
for \((\kappa+2)\)-strong unfoldability, and let
\(\RR_3\) be the corresponding local lottery preparation.
If \(G\subseteq\RR_3\) is generic and \(U=V[G]\), then
\(U\models2^\kappa=\kappa^+\), and, for every
\(\theta\leq\kappa^+\), the full Hechler iteration
computed in \(U\) satisfies
\[
 \PP_\theta\Vdash_U
       \text{\(\kappa\) is \(\Pi^3_1\)-indescribable}.
\]
\end{proposition}
\begin{proof}
By \cite[Corollary 35]{Johnstone}, with \(m=2\), the
initial hypothesis is equivalent to
\((\kappa+2)\)-strong unfoldability. The local Menas
function is supplied by \cite[Theorem 10(2)]{Johnstone}
at that height. Use the size-\(\kappa\), \(\kappa\)-cc
presentation of its lottery preparation described in
\cite[pp.~1229--1230]{Johnstone}.
By \cite[Corollary 38]{Johnstone}, again with \(m=2\), the preparation
preserves \(\Pi^3_1\)-indescribability and makes it
indestructible by \(<\kappa\)-closed,
\(\kappa\)-proper forcing of size at most
\(\beth_{\kappa+1}\).

The preparation preserves \(\kappa^+\), since it is
\(\kappa\)-cc. Every name for a subset of \(\kappa\)
has a nice representative consisting of at most
\(\kappa\) pairs of ordinals below \(\kappa\) and
preparation conditions. There are at most
\((2^\kappa)^\kappa=\kappa^+\) such names in the
ground universe. The old subsets of \(\kappa\) remain
distinct, so \(U\models2^\kappa=\kappa^+\).
Inaccessibility in \(U\) gives \(\beth_\kappa=\kappa\),
and hence the bound in Corollary 38 is
\(\beth_{\kappa+1}=2^\kappa=\kappa^+\), computed in
\(U\).

Fix \(\theta\leq\kappa^+\). In \(U\),
Lemma~\ref{lem:iteration-size} gives a
\(<\kappa\)-closed dense presentation of
\(\PP_\theta\) of size at most \(\kappa^+\).
It is \(\kappa^+\)-cc by Lemma~\ref{lem:iteration}:
compatibility between its conditions agrees with
compatibility in the full iteration, because the
presentation is dense. Since
\(\kappa^{<\kappa}=\kappa\),
\cite[Fact 13(4)]{Johnstone} makes it
\(\kappa\)-proper. Corollary 38 therefore applies
to this presentation. Forcing equivalence gives the
stated conclusion for \(\PP_\theta\).
\end{proof}

\subsection{Eliminating amenable \texorpdfstring{\(C\)}{C}-sequences}

To eliminate amenable \(C\)-sequences, we apply the capture
theorem in an intermediate universe containing the given
sequence, then preserve its stationary capture set through
the tail. The following theorem requires weak compactness
at each proper initial stage, independently of how it is
preserved.

\begin{theorem}\label{thm:conditional}
Let \(\kappa\) be weakly compact and let
\(\lambda=\kappa^+\). Form the Hechler iteration
\(\PP_\lambda\) of Definition~\ref{def:iteration}, and suppose
that
\[
 \PP_\theta\Vdash
       \text{\(\kappa\) is weakly compact}
                     \qquad(\theta<\lambda).
\]
Then \(\PP_\lambda\Vdash\kappa\notin\A_\kappa\).
\end{theorem}

\begin{proof}
Let \(G_\lambda\subseteq\PP_\lambda\) be generic over the
ground universe \(V\), and let
\(G_\theta=G_\lambda\cap\PP_\theta\) for
\(\theta\leq\lambda\). Fix a \(C\)-sequence
\(\vec C=\langle C_\beta:\beta<\kappa\rangle\) in
\(V[G_\lambda]\). By Lemma~\ref{lem:localization}, there is
\(\theta<\lambda\) such that \(\vec C\in V[G_\theta]\),
and it is already a \(C\)-sequence there. Put
\(M=V[G_\theta]\). The hypothesis makes \(\kappa\) weakly
compact in \(M\).

Apply Corollary~\ref{cor:capture} inside \(M\) to the
restriction of \(\vec C\) to \(\Inacc^M(\kappa)\).
At coordinate \(\theta\), the iterand is the full order
\(\D_\kappa^M\). If \(d_\theta\) is its generic function,
then \(M[d_\theta]=V[G_{\theta+1}]\) contains a club
\(D\subseteq\kappa\) and a stationary set
\(T\subseteq\Inacc^M(\kappa)\) satisfying
\[
 D\cap\beta\subseteq C_\beta\qquad(\beta\in T).
\]
Since \(\theta+1<\lambda\), the tail forcing from
that stage to \(\lambda\) is
\(<\kappa\)-closed in \(V[G_{\theta+1}]\) by
Lemma~\ref{lem:iteration}, and it preserves the stationarity
of \(T\) by Lemma~\ref{lem:stationarity}. Closedness and
unboundedness of \(D\), and the displayed inclusions, are
absolute to the final extension: they concern fixed sets of
ordinals. Thus \(\operatorname{Cap}(D,\vec C)\) contains
the still stationary set \(T\) in \(V[G_\lambda]\),
so \(\vec C\) is not amenable. Since the sequence was
arbitrary, \(\kappa\notin\A_\kappa\) in the final extension.
\end{proof}

\begin{proof}[Proof of Theorem~\ref{thm:main}]
By Corollary~\ref{cor:ground}, begin in a universe
\(W\models\ZFC+V=L\) containing a
\(\Pi^2_1\)-indescribable cardinal. Fix the least such
cardinal \(\kappa\) in \(W\).
Proposition~\ref{prop:least}, applied inside \(W\), shows
that \(\kappa\) is not ineffable in \(W=L^W\). Apply
Theorem~\ref{thm:preparation} inside \(W\), and let
\(U=W[H]\) be the resulting prepared universe.

In \(U\), form the \(<\kappa\)-support Hechler iteration
of length \(\lambda=(\kappa^+)^U\), and let
\(G_\lambda\subseteq\PP_\lambda\) be generic over \(U\).
For \(\theta\leq\lambda\), let
\(G_\theta=G_\lambda\cap\PP_\theta\).
Corollary~\ref{cor:iteration-preservation} gives
\[
 U[G_\theta]\models
       \text{\(\kappa\) is \(\Pi^2_1\)-indescribable}
                         \qquad(\theta\leq\lambda).
\]
In particular, weak compactness holds at every proper
initial stage, so Theorem~\ref{thm:conditional} applies
over \(U\). The final universe
\(N=U[G_\lambda]\) therefore satisfies
\[
 N\models
       \text{\(\kappa\) is \(\Pi^2_1\)-indescribable and }
       \kappa\notin\A_\kappa.
\]

Set forcing leaves the ordinals and the constructible
universe unchanged. Hence \(L^N=L^W=W\), since
\(W\models V=L\). The choice of \(\kappa\) in \(W\) gives
\[
 L^N\models\text{\(\kappa\) is not ineffable}.
\]
The forcing theorem yields the asserted relative consistency
implication.
\end{proof}


\begin{corollary}
\label{cor:pi3-consistency}
If \(\ZFC+\) ``there is a \(\Pi^3_1\)-indescribable
cardinal'' is consistent, then so is
\[
 \ZFC+\exists\kappa\,
 \bigl[\kappa\text{ is }\Pi^3_1\text{-indescribable}
       \ \land\ \kappa\notin\A_\kappa
       \ \land\ L\models\text{\(\kappa\) is not ineffable}\bigr].
\]
\end{corollary}
\begin{proof}
As for the \(\Pi^2_1\) hypothesis, the consistency
of \(\Pi^3_1\)-indescribability with \(V=L\) is
recorded in \cite[Section 1, p.~1215]{Johnstone}.
Begin in \(W\models\ZFC+V=L\) and take its least
\(\Pi^3_1\)-indescribable cardinal \(\kappa\).
If \(\kappa\) were subtle in \(W\),
\cite[Lemma 5.4]{HLN} would give a totally
indescribable cardinal below it, contradicting
its minimality. Thus \(\kappa\) is not ineffable
in \(W\), and GCH supplies
\(2^\kappa=\kappa^+\).

Apply Proposition~\ref{prop:pi3-preparation} in
\(W\), and then force with the full Hechler
iteration of length \(\kappa^+\) in the prepared
universe \(U\). The proposition preserves
\(\Pi^3_1\)-indescribability at every initial
stage, including the final stage. In particular,
weak compactness holds at every proper initial
stage, so Theorem~\ref{thm:conditional} gives
\(\kappa\notin\A_\kappa\) in the final universe
\(N\). Since \(W\models V=L\) and both forcings
are set forcings, \(L^N=W\). Consequently,
\(L^N\) still regards \(\kappa\) as non-ineffable.
\end{proof}

\appendix
\section{Preservation under the
  \texorpdfstring{\(\Pi^2_1\)}{Pi-2-1} hypothesis}
\label{sec:preservation}

We prove Theorem~\ref{thm:preparation} by adapting the two-step
lifting argument of \cite[pp.~1229--1232]{Johnstone} to the
full Hechler iteration. The canonical presentation below lets
us identify the iteration computed by the target with the
actual order in the prepared universe.

\subsection{Embedding data}

\begin{definition}[\(\kappa\)-models]\label{def:kappa-model}
Throughout the article, \(\ZFCc\) denotes \(\ZFC\) without
Power Set, with full Separation, Replacement, and Collection,
and with the assertion that every set can be well ordered.
For inaccessible \(\kappa\), a \emph{\(\kappa\)-model} is a
transitive set \(M\) such that
\[
 |M|=\kappa,\qquad \kappa\in M,\qquad M\models\ZFCc,
 \qquad {}^{<\kappa}M\subseteq M.
\]
Size and closure are evaluated in the ambient universe.
\end{definition}

Collection is specified because, without Power Set, it need
not follow from Replacement \cite{GHJ}.

For inaccessible \(\kappa\), \(\Pi^2_1\)-indescribability
is equivalent to the following property: for every
\(\kappa\)-model \(M\), there are a transitive set \(N\)
and an elementary embedding \(j:M\to N\) such that
\[
 \crit(j)=\kappa,\qquad j(\kappa)>\kappa+1,
 \qquad V_{\kappa+1}\subseteq N.
\]
Here the rank segment belongs to the ambient universe. This is
\((\kappa+1)\)-strong unfoldability in
\cite[Definition 3 and Corollary 35]{Johnstone}.

The characterization remains valid when \(\kappa\)-models
are required to satisfy Collection.
Every \(A\subseteq\kappa\) can be placed into a model of
Definition~\ref{def:kappa-model} by collapsing a
\(<\kappa\)-closed elementary submodel of a sufficiently
large \(H_\chi\), of size \(\kappa\), containing \(A\)
and all ordinals at most \(\kappa\). For regular \(\chi\),
\(H_\chi\) satisfies Collection. The embedding property for
these models implies the subset-of-\(\kappa\)
characterization in \cite[Fact 4(4)]{Johnstone}. Conversely,
every model of Definition~\ref{def:kappa-model} satisfies
the \(\kappa\)-model requirements of \cite{Johnstone}, so the same
fact supplies the embedding for each of these models.
Together with \cite[Corollary 35]{Johnstone}, with \(m=1\),
this proves the asserted equivalence.

\begin{definition}[Local Menas property]\label{def:menas}
Let \(\kappa\) be inaccessible and let \(\eta\geq\kappa\)
be an ordinal. A partial function \(f\) from \(\kappa\) to
\(\kappa\) has the \emph{\(\eta\)-strong unfoldability Menas
property} if every \(\kappa\)-model \(M\) containing \(f\)
admits an elementary embedding \(j:M\to N\), with \(N\)
transitive, such that
\[
 \crit(j)=\kappa,\qquad j(\kappa)>\eta,
 \qquad V_\eta\subseteq N,
\]
\[
 \kappa\in\dom(j(f)),\qquad
 j(f)(\kappa)\geq(\beth_\eta)^N.
\]
The beth number in the last inequality is computed in \(N\).
This is the local Menas property at height \(\eta\) of
\cite[Section 3, p.~1221]{Johnstone}. The \emph{strong
unfoldability Menas property} requires this at every ordinal
height \(\eta\geq\kappa\). In this appendix we use only the local
property at height \(\kappa+1\).
\end{definition}

We use the local embedding data of
\cite[Fact 5 and Theorem 10(2)]{Johnstone} in the following form.

\begin{proposition}\label{prop:target}
If \(\kappa\) is \(\Pi^2_1\)-indescribable, there is a
partial local Menas function \(f\) from \(\kappa\) to
\(\kappa\) at height \(\kappa+1\) such that, for each
\(\kappa\)-model \(M\) containing \(f\), one can choose
an elementary embedding \(j:M\to N\) satisfying
\[
 \crit(j)=\kappa,\qquad j(\kappa)>\kappa+1,
 \qquad V_{\kappa+1}\subseteq N,
\]
\[
 {}^\kappa N\subseteq N,\qquad |N|=2^\kappa,
 \qquad j(f)(\kappa)>\mu,
 \qquad \mu=(\beth_{\kappa+1})^N=(2^\kappa)^N.
\]
The target regards \(j(\kappa)\) as inaccessible and has
\(V_{j(\kappa)}^N\) as a set.
\end{proposition}
\begin{proof}
Use \cite[Corollary 35 and Fact 5]{Johnstone} at successor
height \(\kappa+1\). Let \(f_0\) be the partial function
supplied by \cite[Theorem 10(2)]{Johnstone}, and put
\(f(\xi)=f_0(\xi)+1\) for \(\xi\in\dom(f_0)\).
Its values remain below \(\kappa\), and every
\(\kappa\)-model containing \(f\) also contains \(f_0\):
take the unique predecessor of each value. As explained in
\cite[Section 3]{Johnstone}, the Menas property of \(f_0\)
can be retained when selecting the targets of Fact 5.
Definition~\ref{def:menas} therefore gives
\[
 j(f)(\kappa)=j(f_0)(\kappa)+1
              >(\beth_{\kappa+1})^N=\mu.
\]
Inaccessibility and external \(<\kappa\)-closure give
\(V_\kappa\subseteq M\): by induction on rank, each
\(x\in V_\kappa\) is the range of a sequence of fewer
than \(\kappa\) elements already in \(M\).
The same argument puts each \(V_\alpha\), for
\(\alpha<\kappa\), into \(M\). Replacement then
forms \(V_\kappa\) in \(M\), and \(M\) regards
\(\kappa\) as inaccessible. Elementarity gives
inaccessibility of \(j(\kappa)\) in \(N\) and
\(j(V_\kappa)=V_{j(\kappa)}^N\in N\).
\end{proof}

\subsection{A presentation determined by
  \texorpdfstring{\(H_{\kappa^+}\)}{H-kappa-plus}}
\label{subsec:canonical}

For the preservation argument, we need conditions and
side-function names of hereditary size at most \(\kappa\).
Fix a strongly inaccessible \(\kappa\) and write
\(K=H_{\kappa^+}\), calculated in the ground universe.
No assumption on \(2^\kappa\) is needed here.
We use the support conventions of Definition~\ref{def:iteration}
and represent each coordinate by a decided stem and a code
for a nice function name.

A \emph{nice function-name code} over a forcing order \(E\)
is a sequence of the fixed length \(\kappa\),
\[
 B=\langle(A_\zeta,v_\zeta):\zeta<\kappa\rangle,
\]
where \(A_\zeta\) is a maximal antichain in \(E\) and
\(v_\zeta:A_\zeta\to\kappa\). It determines the name
\[
 \dot h_B=
 \{(\check{\langle\zeta,v_\zeta(a)\rangle},a):
                   \zeta<\kappa,\ a\in A_\zeta\}
\]
for the graph of a function from \(\kappa\) to \(\kappa\).
The code itself, rather than the expansion of its check
names, occurs in the conditions below.

\begin{definition}[Canonical presentation]\label{def:canonical}
Define \(\EE_\theta\), for \(\theta\leq\kappa^+\),
recursively, starting with \(\EE_0=\{\varnothing\}\).
A condition in \(\EE_\theta\) is a partial function
\(p\in K\) whose domain is a subset of \(\theta\)
of size \(<\kappa\), such that
\(p\restriction\xi\in\EE_\xi\) for every \(\xi<\theta\).
For \(\xi\in\dom(p)\), its entry has the form
\[
 p(\xi)=(s_\xi,B_\xi),
 \qquad s_\xi\in{}^{<\kappa}\kappa,
\]
where \(B_\xi\) is a nice function-name code over
\(\EE_\xi\) and
\[
 p\restriction\xi\Vdash_{\EE_\xi}
                         s_\xi\subseteq\dot h_{B_\xi}.
\]
At an omitted coordinate we use the empty stem and the
constant zero function, without adding that coordinate to
the domain. The empty function is the greatest condition.

For \(p,q\in\EE_\theta\), define \(q\leq p\) if
\(\supp(p)\subseteq\supp(q)\) and, for every
\(\xi\in\supp(p)\),
\[
 s_\xi^p\subseteq s_\xi^q
 \quad\text{and}\quad
 q\restriction\xi\Vdash_{\EE_\xi}
             \dot h_{B_\xi^p}\leq\dot h_{B_\xi^q}
                         \text{ pointwise}.
\]
Every valid entry extends the greatest condition assigned to
an omitted coordinate, so no additional comparison is required
at a new coordinate.
\end{definition}

\begin{lemma}\label{lem:canonical-nice}
Each \(\EE_\theta\) is forcing equivalent to
\(\PP_\theta\), is \(<\kappa\)-closed and
\(\kappa^+\)-Knaster, and has all its conditions in
\(K\). Every nice function-name code over
\(\EE_\theta\), and every family of at most
\(\kappa\) of its conditions, belongs to \(K\).
\end{lemma}

\begin{proof}
We induct on \(\theta\leq\kappa^+\), constructing maps
\(i_\theta:\EE_\theta\to\PP_\theta\) that preserve
and reflect order, preserve incompatibility, and have dense
range. They will also satisfy
\begin{equation}\label{eq:canonical-coherence}
 \supp(i_\theta(p))=\supp(p),\qquad
 i_\theta(p)\restriction\xi
                  =i_\xi(p\restriction\xi)\quad(\xi<\theta).
\end{equation}
For \(\theta=0\), use the identity map. The only condition
is empty, and every nice code has hereditary size at most
\(\kappa\). Suppose now that \(\theta>0\) and that the
lemma and the maps have been established at all earlier
stages.

\emph{Hereditary size.}
Write \(T(x)=\tc(\{x\})\). If \(\xi<\theta\) and
\(B=\langle(A_\zeta,v_\zeta):\zeta<\kappa\rangle\)
is a nice code over \(\EE_\xi\), the induction
hypothesis gives
\[
 |A_\zeta|\leq\kappa,\qquad
 |T(a)|\leq\kappa\quad(a\in A_\zeta).
\]
The ordinal indices and values contribute at most
\(\kappa\) elements, and coding sequences and ordered
pairs adds only finitely many components per entry. Hence
\begin{equation}\label{eq:canonical-code-size}
 |T(B)|\leq
 \kappa+\sum_{\zeta<\kappa}\sum_{a\in A_\zeta}
                         (\aleph_0+|T(a)|)
 \leq\kappa.
\end{equation}
For a partial function \(p\) with domain
\(u\in[\theta]^{<\kappa}\) and entries of the required
form, each \(\xi\in u\) is below \(\kappa^+\),
so \(|T(\xi)|\leq\kappa\). Its stems have hereditary
size less than \(\kappa\). Thus
\begin{equation}\label{eq:canonical-condition-size}
 \begin{split}
 |T(p)|&\leq\aleph_0+\sum_{\xi\in u}
   \bigl(\aleph_0+|T(\xi)|+|T(s_\xi)|+|T(B_\xi)|\bigr)\\
 &\leq\aleph_0+|u|\cdot\kappa\leq\kappa.
 \end{split}
\end{equation}
In particular, the requirement \(p\in K\) excludes no
partial function with the prescribed entries and support
size. These estimates apply to the conditions constructed
below.

\emph{The map and its order.}
At every earlier stage \(\xi\), translate
\(\EE_\xi\)-names recursively by
\[
 \tau^{i_\xi}=
 \{(\rho^{i_\xi},i_\xi(a)):(\rho,a)\in\tau\}.
\]
Forcing equivalence at that stage gives
\begin{equation}\label{eq:canonical-forcing-translation}
 r\Vdash_{\EE_\xi}\varphi(\tau)
 \quad\Longleftrightarrow\quad
 i_\xi(r)\Vdash_{\PP_\xi}\varphi(\tau^{i_\xi}).
\end{equation}
For a nice code, the translated graph is, up to forced
equality of check names,
\[
 \dot h_B^{i_\xi}=
 \{(\check{\langle\zeta,v_\zeta(a)\rangle},i_\xi(a)):
                         \zeta<\kappa,\ a\in A_\zeta\}.
\]

For \(p\in\EE_\theta\), give \(i_\theta(p)\) the
same domain as \(p\). At a coordinate \(\xi\) in
this domain, validity and
\eqref{eq:canonical-forcing-translation} show that
\[
 i_\xi(p\restriction\xi)\Vdash_{\PP_\xi}
  \langle\check s_\xi^p,\dot h_{B_\xi^p}^{i_\xi}\rangle
                                      \in\dot{\mathbb Q}_\xi.
\]
Fullness of the iterand name therefore supplies a
representative \(i_\theta(p)(\xi)\) in its domain with
\begin{equation}\label{eq:canonical-map-coordinate}
 i_\xi(p\restriction\xi)\Vdash_{\PP_\xi}
 i_\theta(p)(\xi)=
       \langle\check s_\xi^p,\dot h_{B_\xi^p}^{i_\xi}\rangle.
\end{equation}
Choose this representative once for each triple
\((p\restriction\xi,s_\xi^p,B_\xi^p)\). This makes
the choices coherent across stages and gives
\eqref{eq:canonical-coherence}. Coordinate validity and
the support bound give \(i_\theta(p)\in\PP_\theta\).

If \(q\leq p\) in \(\EE_\theta\), induction gives
\(i_\xi(q\restriction\xi)\leq
i_\xi(p\restriction\xi)\) for every \(\xi<\theta\).
This initial condition forces both coordinate identities
in \eqref{eq:canonical-map-coordinate}. The stem inclusion
and side-function inequality defining \(q\leq p\) then
give the coordinate inequalities for
\(i_\theta(q)\leq i_\theta(p)\). Conversely, if the
latter inequality holds, restriction and order reflection
at earlier stages give
\(q\restriction\xi\leq p\restriction\xi\).
Substituting the coordinate identities into the forced
Hechler inequality yields stem inclusion and, by
\eqref{eq:canonical-forcing-translation}, the required
side-function inequality in \(\EE_\xi\). The domains
are unchanged. Thus
\begin{equation}\label{eq:canonical-order-reflection}
 i_\theta(q)\leq_{\PP_\theta}i_\theta(p)
              \quad\Longleftrightarrow\quad
 q\leq_{\EE_\theta}p.
\end{equation}

\emph{Density and forcing equivalence.}
Given \(p\in\PP_\theta\), choose
\(p'\in\mathcal E_\theta\) below \(p\), using
Lemma~\ref{lem:iteration}, and put \(u=\supp(p')\).
For \(\xi\in u\), let \(s_\xi\) be its decided
stem. Extend the side-function name to a globally total
\(\PP_\xi\)-name \(\dot f_\xi\): on a maximal
antichain containing \(p'\restriction\xi\), retain
that name at \(p'\restriction\xi\) and use zero at
all other members. Then
\[
 \mathbf1_\xi\Vdash\dot f_\xi:\kappa\to\kappa,
 \qquad p'\restriction\xi\Vdash
          p'(\xi)=\langle\check s_\xi,\dot f_\xi\rangle.
\]

For every \(\zeta<\kappa\), conditions \(a\in\EE_\xi\)
whose images decide \(\dot f_\xi(\zeta)\) are dense
in \(\EE_\xi\). Indeed, below \(i_\xi(r)\),
first decide this value in \(\PP_\xi\), then use the
dense range to choose an image below the deciding
condition; order reflection puts its preimage below \(r\).
Choose a maximal antichain \(A_\zeta\) of such
conditions and record the decided values in
\(v_\zeta:A_\zeta\to\kappa\). Since \(i_\xi\)
preserves incompatibility and has dense range,
\(i_\xi[A_\zeta]\) is a maximal antichain in
\(\PP_\xi\). More explicitly, for any
\(t\in\PP_\xi\), choose \(i_\xi(w)\leq t\);
maximality in \(\EE_\xi\) supplies a common extension
of \(w\) and some \(a\in A_\zeta\), whose image
witnesses compatibility with \(t\). Therefore the code
\(B_\xi=\langle(A_\zeta,v_\zeta):\zeta<\kappa\rangle\)
satisfies
\[
 \mathbf1_\xi\Vdash_{\PP_\xi}
                    \dot h_{B_\xi}^{i_\xi}=\dot f_\xi.
\]
The earlier chain condition and
\eqref{eq:canonical-code-size} give \(B_\xi\in K\).

Define \(q\) on \(u\) by \(q(\xi)=(s_\xi,B_\xi)\).
Inductively, suppose
\(i_\xi(q\restriction\xi)\leq p'\restriction\xi\).
For \(\xi\in u\), this condition forces
\(s_\xi\subseteq\dot h_{B_\xi}^{i_\xi}\), so
\eqref{eq:canonical-forcing-translation} verifies the
entry's validity in \(\EE_{\xi+1}\). It also forces
the new coordinate of the image to equal \(p'(\xi)\),
by \eqref{eq:canonical-map-coordinate}. Hence
\[
 i_{\xi+1}(q\restriction(\xi+1))
                       \leq p'\restriction(\xi+1).
\]
At omitted coordinates this inequality is unchanged.
At a limit, all preceding validity tests hold and the
domain remains a subset of \(u\), of size \(<\kappa\).
The hereditary-size estimate consequently gives
\[
 q\in\EE_\theta,\qquad
 \supp(q)=u,\qquad i_\theta(q)\leq p'\leq p.
\]
Thus the range of \(i_\theta\) is dense.

If two images have a common extension, density gives an
image below it, and order reflection gives a common
extension of their preimages. The converse follows from
order preservation. In particular, \(i_\theta\)
preserves incompatibility. The associated generic filters
correspond by
\[
 H=\{r\in\PP_\theta:(\exists p\in G)\ i_\theta(p)\leq r\},
 \qquad G=\{p\in\EE_\theta:i_\theta(p)\in H\}.
\]
For a generic \(H\), there is an image in \(H\) below
every member of \(H\): the images below a given condition,
together with the conditions incompatible with it, form a
dense set. This fact and order reflection give directedness
of the second filter and show that the assignments are
inverse. The image of a dense subset of \(\EE_\theta\)
is dense in \(\PP_\theta\); for a dense
\(D\subseteq\PP_\theta\), the set of \(p\in\EE_\theta\)
whose images extend a member of \(D\) is dense. Thus both
filters are generic. Induction on name rank gives
\(\tau^G=(\tau^{i_\theta})^H\). This proves forcing
equivalence and establishes
\eqref{eq:canonical-forcing-translation} at stage
\(\theta\).

Apply the \(\kappa^+\)-Knaster property of
\(\PP_\theta\), from Lemma~\ref{lem:iteration}, to
the images of a \(\kappa^+\)-sequence in \(\EE_\theta\).
Compatibility of images is equivalent to compatibility
of preimages, so \(\EE_\theta\) is
\(\kappa^+\)-Knaster. Its antichains therefore have
size at most \(\kappa\), and
\eqref{eq:canonical-code-size} applies also to nice
codes over \(\EE_\theta\). If
\(X\subseteq\EE_\theta\) and \(|X|\leq\kappa\), then
\[
 |T(X)|\leq\aleph_0+\sum_{p\in X}(1+|T(p)|)
                           \leq\kappa.
\]
This proves the remaining hereditary-size assertions.

\emph{Closure.}
Let \(\langle p_i:i<\delta\rangle\) be decreasing in
\(\EE_\theta\), where \(0<\delta<\kappa\), and put
\[
 u=\bigcup_{i<\delta}\supp(p_i).
\]
Regularity of \(\kappa\) gives \(|u|<\kappa\).
For \(\xi\in u\), denote the coordinate stems and
side-function names by \(s_\xi^i\) and
\(\dot h_\xi^i\), using the empty stem and zero
function when \(\xi\) is omitted. The stems are
increasing, so
\[
 t_\xi=\bigcup_{i<\delta}s_\xi^i\in{}^{<\kappa}\kappa.
\]
The earlier \(<\kappa\)-closure preserves regularity
of \(\kappa\). As each side-function name is globally
total, we may choose a name \(\dot g_\xi\) with
\[
 \mathbf1_{\EE_\xi}\Vdash
 \dot g_\xi(\zeta)=\sup_{i<\delta}\dot h_\xi^i(\zeta)<\kappa
                                      \quad(\zeta<\kappa).
\]
Choose maximal antichains deciding these values and let
\(C_\xi\) be the resulting nice code. The earlier
chain condition gives
\[
 C_\xi\in K,\qquad
 \mathbf1_{\EE_\xi}\Vdash\dot h_{C_\xi}=\dot g_\xi.
\]

Define \(q\) on \(u\) by \(q(\xi)=(t_\xi,C_\xi)\).
Suppose inductively that
\(q\restriction\xi\leq p_i\restriction\xi\) for
all \(i<\delta\). Below this initial condition, the
side functions increase with \(i\), and each extends
its stem. If \(\zeta<\dom(t_\xi)\), choose
\(i_0<\delta\) with \(\zeta<\dom(s_\xi^{i_0})\).
For every \(j\geq i_0\), stem extension and
normalization give
\(\dot h_\xi^j(\zeta)=t_\xi(\zeta)\); for
\(j<i_0\), the forced side-function inequality gives
\(\dot h_\xi^j(\zeta)\leq t_\xi(\zeta)\).
Consequently
\[
 q\restriction\xi\Vdash
 \dot h_{C_\xi}(\zeta)
      =\sup_{j<\delta}\dot h_\xi^j(\zeta)=t_\xi(\zeta).
\]
This verifies normalization of the new entry. Since
\(t_\xi\) extends every \(s_\xi^i\) and
\(\dot h_{C_\xi}\) pointwise dominates every
\(\dot h_\xi^i\), the same initial condition verifies
\[
 q\restriction(\xi+1)\leq p_i\restriction(\xi+1)
                                      \quad(i<\delta).
\]
At limits, the domain has size \(<\kappa\) and all
earlier validity tests remain satisfied. By
\eqref{eq:canonical-condition-size}, \(q\in K\),
so \(q\in\EE_\theta\) and \(q\leq p_i\) for
all \(i<\delta\). The empty condition bounds an
empty sequence. This proves closure and completes the
induction.
\end{proof}

\begin{lemma}
\label{lem:canonical-absolute}
Suppose \(W\subseteq U\) is a transitive model of
\(\ZFCc\), with the correct \(\kappa^+\), and
\[
 H_{\kappa^+}^W=H_{\kappa^+}^U=K\in W.
\]
Then
\[
 (\EE_\theta)^W=(\EE_\theta)^U
                 \qquad(\theta\leq\kappa^+),
\]
including the order relations.
\end{lemma}
\begin{proof}
We describe the recursion using tests on sets that are
common to the two models. All candidate conditions lie
in \(K\), as do their restrictions, stems, and codes.
For such a candidate, the assertion that its support has
size \(<\kappa\) is absolute. An external bijection
from some \(\delta<\kappa\) onto that support has
hereditary size at most \(\kappa\), hence belongs to
\(K\subseteq W\); conversely, an internal bijection
remains a bijection in \(U\).

Suppose the earlier conditions and orders are common.
Compatibility in \(\EE_\xi\) is then absolute, since
it asks for a common extension in that same set. The
statement that a given \(A\) is a maximal antichain
also quantifies over the common set \(\EE_\xi\).
Thus the nice-code tests agree. The forcing assertions
in Definition~\ref{def:canonical} can be expressed
directly in terms of these codes. For example,
\[
 r\Vdash\dot h_B(\zeta)=\check\gamma
 \quad\Longleftrightarrow\quad
 (\forall a\in A_\zeta)
   \bigl(\Comp(r,a)\Rightarrow v_\zeta(a)=\gamma\bigr).
\]
Maximality ensures that the compatible members decide
all possible values below \(r\). Similarly,
\(r\Vdash\dot h_B(\zeta)\leq\dot h_C(\zeta)\)
holds exactly when there are no
\(a\in A_\zeta^B\), \(b\in A_\zeta^C\) with
\(v_\zeta^B(a)>v_\zeta^C(b)\) and a common
extension of \(r,a,b\). These tests quantify over
the common earlier order. Testing equality at each
position in a stem gives normalization; testing the
inequality at every \(\zeta<\kappa\) gives the
pointwise comparison. Stem inclusion is absolute as well.

At stage \(\theta\), Separation on \(K\) gives
the conditions satisfying these tests, and Separation
on \(K\times K\) gives their order. These fixed sets
belong to \(W\), and the tests define a unique next stage
from each earlier sequence. Transfinite recursion using
Replacement therefore collects the stages through
\(\kappa^+\) inside \(W\). Here \(\kappa^+\)
is a parameter in \(W\), while all coordinates
appearing in conditions are below \(\kappa^+\).
Induction gives equality at every stage, including
limits: the support-size test places no boundedness
requirement on the domain.
\end{proof}

\subsection{Agreement with the prepared universe}

\begin{lemma}\label{lem:target-H}
Suppose \(N\) is transitive and satisfies \(\ZFCc\),
\(V_{\kappa+1}\subseteq N\), and \(N\) has an
internally inaccessible \(\nu>\kappa\) with
\(V_\nu^N\in N\). Then
\[
 (\kappa^+)^N=\kappa^+,
 \qquad H_{\kappa^+}^N=H_{\kappa^+},
 \qquad H_{\kappa^+}\in N.
\]
\end{lemma}
\begin{proof}
The rank segment \(V_\nu^N\) supplies the internal sets
\(\mathcal P(\kappa)^N\), \((\kappa^+)^N\), and
\(H_{(\kappa^+)^N}^N\). Since \(V_{\kappa+1}\subseteq N\),
every actual subset of \(\kappa\) belongs to \(N\).
Fix in \(N\) a bijection from \(\kappa\times\kappa\) to
\(\kappa\). Every well-order of a subset of \(\kappa\) is
therefore coded by a subset of \(\kappa\) in \(N\), and its
Mostowski collapse in \(N\) is its actual collapse. Thus every
\(\alpha<\kappa^+\) belongs to \(N\) and has an internal
enumeration of length at most \(\kappa\). It follows that
\((\kappa^+)^N\geq\kappa^+\). Strict inequality would give
an internal enumeration of the actual \(\kappa^+\) by
\(\kappa\), also an enumeration in the ambient universe.
Hence \((\kappa^+)^N=\kappa^+\).

For \(x\in H_{\kappa^+}\), code \(\tc(\{x\})\) by a
well-founded extensional relation on a subset of \(\kappa\),
with a distinguished element corresponding to \(x\). This
code belongs to \(N\). Its internal collapse gives \(x\)
and an enumeration witnessing its hereditary size at most
\(\kappa\). Conversely, any such internal enumeration is an
actual enumeration. Therefore
\(H_{\kappa^+}^N=H_{\kappa^+}\); the common set belongs
to \(N\) because the internal set exists.
\end{proof}

\begin{lemma}\label{lem:target-afterprep}
Let \(j:M\to N\) be as in
Proposition~\ref{prop:target}, and let
\(G\subseteq\RR\) be generic over \(V\). Put
\(U=V[G]\). Then
\[
 H_{\kappa^+}^{N[G]}=H_{\kappa^+}^{U}\in N[G].
\]
If \(V\models2^\kappa=\kappa^+\), then
\(U\models2^\kappa=\kappa^+\).
\end{lemma}
\begin{proof}
The preparation below \(\kappa\) is the same in \(N\)
and \(V\), so \(G\) is also \(N\)-generic. Every relation
on \(\kappa\) in \(U\) has a nice \(\RR\)-name with at
most \(\kappa\) pairs. The conditions in this name belong
to \(V_\kappa\), so the name has hereditary size at most
\(\kappa\). By Lemma~\ref{lem:target-H}, it belongs to
\(N\). Consequently \(N[G]\) contains every relation on
\(\kappa\) in \(U\).

The \(\kappa\)-cc of \(\RR\) preserves \(\kappa^+\)
both in \(V\) and internally in \(N\). Internally,
\(\RR\) has size \(\kappa<j(\kappa)\), so
\(j(\kappa)\) remains inaccessible in \(N[G]\), and
\(V_{j(\kappa)}^{N[G]}\in N[G]\). The coding and collapse
argument of Lemma~\ref{lem:target-H} now proves the
assertion about \(H_{\kappa^+}\).

Finally, if \(2^\kappa=\kappa^+\) in \(V\), the number
of nice names for subsets of \(\kappa\) is at most
\((2^\kappa)^\kappa=\kappa^+\). Together with preservation
of \(\kappa^+\) and the old subsets of \(\kappa\), this
gives \(U\models2^\kappa=\kappa^+\).
\end{proof}

The proof uses the external size \(|N|^V=\kappa^+\) to
enumerate dense sets. The lottery bound, however, is the
internal cardinal \(\mu=(2^\kappa)^N\). We do not assume
that \(N\models2^\kappa=\kappa^+\).

\subsection{Lifting the embedding}

\begin{proof}[Proof of Theorem~\ref{thm:preparation}]
Fix \(\theta\leq\kappa^+\). In the prepared universe
use \(Q=\EE_\theta\), forcing equivalent to
\(\PP_\theta\) by Lemma~\ref{lem:canonical-nice}, and
let \(\dot Q\) be its canonical \(\RR\)-name. In \(V\),
fix an \(\RR*\dot Q\)-name \(\dot A\) forced to be a
subset of \(\kappa\), and an arbitrary condition \(r\).
Choose
\[
 X\prec H_\Lambda,
 \quad |X|=\kappa,
 \quad {}^{<\kappa}X\subseteq X,
 \quad \kappa+1\subseteq X,
\]
with \(\kappa,f,\RR,\dot Q,\dot A,r\in X\), where
\(\Lambda\) is sufficiently large for the forcing
constructions under consideration. Since \(\RR\) is
\(\kappa\)-cc and forces that \(Q\) is \(\kappa^+\)-cc,
\(\RR*\dot Q\) is \(\kappa^+\)-cc. Every maximal
antichain in \(X\) consequently has size at most
\(\kappa\) and is contained in \(X\): an enumeration
by \(\kappa\) can be chosen in \(X\), and
\(\kappa\subseteq X\). Thus the restriction to \(X\)
of any actual generic filter \(G*g\) containing \(r\)
is \(X\)-generic. Here \(G\) is \(V\)-generic for
\(\RR\), and \(g\) is \(U\)-generic for the full
order \(Q\), where \(U=V[G]\).

Let \(\pi:X\to M\) be the transitive collapse. The
choice of \(X\) makes \(M\) a \(\kappa\)-model. Since
\(V_\kappa\subseteq X\), the map \(\pi\) fixes
\(\kappa,f,\RR\) and every condition of \(\RR\). Put
\[
 \dot Q_0=\pi(\dot Q),\qquad
 \dot A_0=\pi(\dot A).
\]
By \cite[Lemma 19]{Johnstone}, the image of
\((G*g)\cap X\) is an \(M\)-generic filter
\(G*g_0\) on \(\RR*\dot Q_0\), and
\[
 A=\dot A^{G*g}=\dot A_0^{G*g_0}
                                      \in M[G][g_0].
\]
Writing \(Q_0=\dot Q_0^G\), we have
\(M[G]\models\text{\(Q_0\) is }<\kappa\text{-closed}\).
This collapsed order may differ from \(Q\); the
filter \(g_0\) will be used in the domain, while the
actual filter \(g\) will be used in the target.

Choose \(j:M\to N\) as in
Proposition~\ref{prop:target}. Lemmas
\ref{lem:target-afterprep} and
\ref{lem:canonical-absolute} give
\[
 Q=(\EE_\theta)^U=(\EE_\theta)^{N[G]}\in N[G].
\]
Indeed, \(N[G]\) has the correct \(\kappa^+\) and
contains the common set \(H_{\kappa^+}\) as an element.
In particular, \(\kappa^+<j(\kappa)\).

We next verify the internal size bound needed for the
lottery. The size-\(\kappa\), \(\kappa\)-cc preparation
preserves the internal cardinal
\(\mu=(2^\kappa)^N\) and its value as \(2^\kappa\),
by the nice-name count in
Lemma~\ref{lem:target-afterprep}. Internally in \(N[G]\),
\(H_{\kappa^+}\) has cardinality \(2^\kappa=\mu\):
its elements are coded by relations on \(\kappa\),
and it contains all subsets of \(\kappa\). Every
condition of \(Q\) belongs to this transitive set,
and the order is a relation on its underlying set.
Hence
\[
 N[G]\models |\tc(\{Q\})|\leq\mu
                   \quad\text{and}\quad
                   Q\text{ is }<\kappa\text{-closed}.
\]
Since \(j(f)(\kappa)>\mu\), the actual order \(Q\)
belongs to the stage-\(\kappa\) lottery of \(j(\RR)\).
This stage is eligible because
\(\kappa\in\dom(j(f))\) and
\(j(f)''\kappa=f''\kappa\subseteq\kappa\).
Let \(\dot Q_N\) be the canonical \(\RR\)-name for
\(\EE_\theta\) formed in \(N\). The internal size
and closure calculations hold in every
\(\RR\)-generic extension of \(N\), so the greatest
condition forces that \(\dot Q_N\) is eligible for
this lottery. Choose \(z\) with trivial restriction
below \(\kappa\) and with stage-\(\kappa\) coordinate
choosing \(\dot Q_N\). Since \(\dot Q_N^G=Q\), below
\(z\) the image preparation factors as
\[
 j(\RR)\restriction z
       \cong\RR*\dot Q_N*\RR_{\mathrm{tail}}.
\]
Every later eligible stage lies above
\(j(f)(\kappa)\). The closure of the iterands and
Easton support therefore make the tail internally
closed under descending sequences of length at most
\(\mu\), as in \cite[p.~1229]{Johnstone}.

\emph{First lift.} Work in \(T=V[G][g]\). Apply
\cite[Fact 14(3)]{Johnstone} twice, with
\(\delta=\kappa^+\). Initially
\({}^\kappa N\subseteq N\); the two forcings are
\(\kappa^+\)-cc, and their conditions belong to the
successive targets:
\(\RR\subseteq N\) and
\(Q\subseteq H_{\kappa^+}^U\subseteq N[G]\).
Thus
\[
 {}^\kappa(N[G][g])\subseteq N[G][g]
                                           \quad\text{in }T.
\]
Evaluation of the names in \(N\) gives
\(|N[G][g]|^T\leq\kappa^+\). In fact equality holds,
since \(N\subseteq N[G][g]\) has size \(\kappa^+\)
and the two forcings preserve \(\kappa^+\).

List the dense subsets of \(\RR_{\mathrm{tail}}\) that belong
to \(N[G][g]\) in a sequence of length at most \(\kappa^+\),
and build a descending chain meeting them. At a limit
\(0<\alpha<\kappa^+\), choose in the target a
surjection from \(\kappa\) onto \(\alpha\). Such a
surjection exists because the target has the correct
\(\kappa^+\). The preceding conditions, listed along
that surjection, form a \(\kappa\)-sequence in the
target by external closure. Using the surjection,
the target reconstructs the original
\(\alpha\)-sequence. Since \(\alpha<\kappa^+\leq\mu\),
internal tail closure provides a lower bound. No
lower bound for the entire \(\kappa^+\)-sequence is
needed. The chain generates an
\(N[G][g]\)-generic filter \(G_{\mathrm{tail}}\) in
\(T\).

Under the displayed factorization, let
\(j(G)=G*g*G_{\mathrm{tail}}\). This is an
\(N\)-generic filter on \(j(\RR)\) containing
\(j''G=G\), so the lifting criterion gives
\[
 j:M[G]\longrightarrow N[j(G)].
\]
The target \(N[j(G)]\) still has external size at
most \(\kappa^+\). Since \(G_{\mathrm{tail}}\in T\),
\cite[Fact 14(1)]{Johnstone}, applied in \(T\), also
preserves its external \(\kappa\)-sequence closure.

\emph{Second lift.} The \(\kappa\)-cc preparation
preserves the external \(<\kappa\)-closure of
\(M[G]\), by \cite[Fact 14(3)]{Johnstone} with
\(\delta=\kappa\). The actual order \(Q\) is
\(<\kappa\)-distributive, so \(M[G]\) remains
\(<\kappa\)-closed in \(T\). Finally,
\cite[Fact 14(1)]{Johnstone}, applied in \(T\) to
the \(M[G]\)-generic filter \(g_0\), shows that
\(M[G][g_0]\) is externally \(<\kappa\)-closed.
Since \(|M|=\kappa\), evaluating the names in \(M\)
gives \(|M[G][g_0]|\leq\kappa\). The reverse inequality
follows from \(\kappa\subseteq M[G][g_0]\), so this
model has size \(\kappa\) and is a \(\kappa\)-model
in \(T\).

We need a descending cofinal chain in \(g_0\).
First, every descending sequence of fewer than
\(\kappa\) members of \(g_0\) has a lower bound
in \(g_0\). To prove this, external closure places
the sequence in \(M[G][g_0]\). Internal
\(<\kappa\)-distributivity of \(Q_0\) then places
it in \(M[G]\), since its entries belong to the
old set \(Q_0\). For each entry \(q\), the set
\[
 \{p\in Q_0:p\leq q\text{ or }p\perp q\}
\]
is dense open. Their intersection is dense in
\(M[G]\) by distributivity. A member of its
intersection with \(g_0\) is compatible with every
entry, and hence extends every entry.

The filter \(g_0\) is directed and has external
size at most \(\kappa\). Enumerating its members
and using the preceding closure at limit stages,
choose a descending cofinal chain
\(\langle q_\xi:\xi<\kappa\rangle\) in \(g_0\).
External \(\kappa\)-closure of \(N[j(G)]\) puts
\(\langle j(q_\xi):\xi<\kappa\rangle\) in that
target. By elementarity,
\(N[j(G)]\models\text{\(j(Q_0)\) is }<j(\kappa)\text{-closed}\).
Since \(\kappa<j(\kappa)\), there is a condition
\(q^*\in j(Q_0)\) below the entire image chain,
and therefore below every member of \(j''g_0\).

Starting below \(q^*\), diagonalize through the
dense subsets of \(j(Q_0)\) in \(N[j(G)]\).
There are at most \(\kappa^+\) such sets. The
target still has the correct \(\kappa^+\): its
old enumerations of ordinals below \(\kappa^+\)
remain, and any enumeration of the actual
\(\kappa^+\) by \(\kappa\) in the transitive target
would also be one in \(T\). At each
limit below \(\kappa^+\), the same sequence-coding
argument as in the first lift puts the preceding
chain in the target; its internal
\(<j(\kappa)\)-closure supplies a lower bound,
since \(\kappa^+<j(\kappa)\). This produces in
\(T\) an \(N[j(G)]\)-generic filter \(h\) with
\(q^*\in h\). Consequently \(j''g_0\subseteq h\),
and the embedding lifts to
\[
 j^*:M[G][g_0]\longrightarrow N[j(G)][h].
\]

\emph{Strength of the lifted embedding.} Since
\(\kappa\) is inaccessible in \(U\) and \(Q\) is
\(<\kappa\)-closed, \(V_\kappa^T=V_\kappa^U\).
Fix in \(U\) an enumeration of \(V_\kappa^U\)
of length \(\kappa\). Each subset of
\(V_\kappa^U\) in \(T\) has a nice \(Q\)-name
with at most \(\kappa\) pairs: choose one
antichain deciding membership for each entry of
the enumeration, and use the \(\kappa^+\)-cc.
Every condition in the name has hereditary size
at most \(\kappa\), by
Lemma~\ref{lem:canonical-nice}. The enumeration,
the check names for its entries, and the nice
name itself therefore have hereditary size at
most \(\kappa\). They belong to the common set
\(H_{\kappa^+}^U=H_{\kappa^+}^{N[G]}\), and their
evaluation by the actual filter \(g\) gives
\[
 V_{\kappa+1}^T\subseteq N[G][g]
                         \subseteq N[j(G)][h].
\]
Moreover \(j^*(\kappa)=j(\kappa)>\kappa+1\),
and \(A\) belongs to the transitive domain of
size \(\kappa\). Since \(\dot A\) was arbitrary,
\cite[Fact 4(5) and Corollary 35]{Johnstone}
give \(\Pi^2_1\)-indescribability in \(T\).
The construction works for every generic filter
containing the arbitrary condition \(r\), and
for every \(\theta\leq\kappa^+\), proving the theorem.
\end{proof}


\begin{remark}\label{rem:remove-gch}
The initial equality \(2^\kappa=\kappa^+\) can be
arranged by set forcing. Present
\(\operatorname{Coll}(\kappa^+,2^\kappa)\) as the
order of partial functions from \(\kappa^+\) to
the old \(2^\kappa\) of size \(<\kappa^+\), ordered
by extension. This order is \(<\kappa^+\)-closed;
the union of the generic filter is a surjection
onto the old \(2^\kappa\). It preserves
\(\kappa^+\) and adds no ordinal sequences of
length at most \(\kappa\), hence no subsets of
\(\kappa\). Thus \(2^\kappa=\kappa^+\) in the
extension.

To verify preservation of indescribability,
note that \(V_{\kappa+1}\) and all relation
codes on \(\kappa\) are unchanged. Any
\(\kappa\)-model in the extension has a code
by a well-founded extensional relation on
\(\kappa\). This code is old, and its absolute
collapse is a ground-model \(\kappa\)-model:
the old code already witnesses its size, its
theory is absolute, and its closure in the
extension includes all old short sequences.
An old \((\kappa+1)\)-strong unfoldability
embedding for that model still has the required
rank inclusion and is still elementary. Hence
\(\Pi^2_1\)-indescribability is preserved; compare
\cite[Lemma 6]{Johnstone}. One may then apply
Theorem~\ref{thm:preparation}. Our consistency
proof starts in \(L\), where this preliminary
collapse is unnecessary.
\end{remark}

\section*{Acknowledgments}
I thank Assaf Rinot for discussions on this problem during my
postdoctoral fellowship under his supervision at Bar-Ilan
University (2019--2021).

\section*{Declaration of generative AI and AI-assisted technologies
in the manuscript preparation process}
OpenAI's Codex, using the Sol 6.1 model, was used to assist
with literature searches, exploration of proof strategies,
development of mathematical arguments, and checks of proofs,
as well as the organization, English-language revision,
bibliography, and LaTeX preparation of this manuscript.
The author takes full responsibility for the results and the
accuracy and integrity of the article.

\begin{thebibliography}{99}
\raggedright

\bibitem{BNZRad}
Omer Ben-Neria and Jing Zhang,
\emph{Compactness and guessing principles in the Radin extensions},
Journal of Mathematical Logic \textbf{23}(2) (2023), Article 2250024.
\href{https://doi.org/10.1142/S0219061322500246}
     {doi:10.1142/S0219061322500246}.
The question number cited here refers to
\href{https://arxiv.org/abs/2105.01037v2}{arXiv:2105.01037v2}.

\bibitem{BNZApprox}
Omer Ben-Neria and Jing Zhang,
\emph{Approximating diamond principles on products at an
inaccessible cardinal},
Transactions of the American Mathematical Society
\textbf{376}(8) (2023), 5923--5948.
\href{https://doi.org/10.1090/tran/8945}
     {doi:10.1090/tran/8945}.

\bibitem{BBFM}
J\"org Brendle, Andrew Brooke-Taylor, Sy-David Friedman,
and Diana Montoya,
\emph{Cicho\'n's diagram for uncountable cardinals},
Israel Journal of Mathematics \textbf{225}(2) (2018), 959--1010.
\href{https://doi.org/10.1007/s11856-018-1688-y}
     {doi:10.1007/s11856-018-1688-y}.
The definition numbers cited here refer to the
\href{https://eprints.whiterose.ac.uk/id/eprint/118463/1/cichon-large(revised).pdf}
     {accepted manuscript}.

\bibitem{BR}
Ari Meir Brodsky and Assaf Rinot,
\emph{Distributive Aronszajn trees},
Fundamenta Mathematicae \textbf{245}(3) (2019), 217--291.
\href{https://doi.org/10.4064/fm542-4-2018}
     {doi:10.4064/fm542-4-2018}.

\bibitem{Cummings}
James Cummings,
\emph{Iterated forcing and elementary embeddings},
in: Matthew Foreman and Akihiro Kanamori (eds.),
\emph{Handbook of Set Theory}, Vol.~2,
Springer, Dordrecht, 2010, pp.~775--883.
\href{https://doi.org/10.1007/978-1-4020-5764-9_13}
     {doi:10.1007/978-1-4020-5764-9\_13}.

\bibitem{GHJ}
Victoria Gitman, Joel David Hamkins, and Thomas A. Johnstone,
\emph{What is the theory ZFC without power set?},
Mathematical Logic Quarterly \textbf{62}(4--5) (2016), 391--406.
\href{https://doi.org/10.1002/malq.201500019}
     {doi:10.1002/malq.201500019}.

\bibitem{HLN}
Peter Holy, Philipp L\"ucke, and Ana Njegomir,
\emph{Small embedding characterizations for large cardinals},
Annals of Pure and Applied Logic \textbf{170}(2) (2019), 251--271.
\href{https://doi.org/10.1016/j.apal.2018.10.002}
     {doi:10.1016/j.apal.2018.10.002}.
Author's final manuscript:
\url{https://www.dmg.tuwien.ac.at/holy/SmallEmbChar.pdf}.

\bibitem{IR}
Tanmay Inamdar and Assaf Rinot,
\emph{Was Ulam right? I: Basic theory and subnormal ideals},
Topology and its Applications \textbf{323} (2023), Article 108287.
\href{https://doi.org/10.1016/j.topol.2022.108287}
     {doi:10.1016/j.topol.2022.108287}.
The definition and conjecture numbers cited here refer to
arXiv:2107.05136v2,
\url{https://arxiv.org/abs/2107.05136v2}.

\bibitem{Johnstone}
Thomas A. Johnstone,
\emph{Strongly unfoldable cardinals made indestructible},
The Journal of Symbolic Logic \textbf{73}(4) (2008), 1215--1248.
\href{https://doi.org/10.2178/jsl/1230396915}
     {doi:10.2178/jsl/1230396915}.

\bibitem{LL}
Paul Larson and Chris Lambie-Hanson,
\emph{Strongly increasing sequences},
The Journal of Symbolic Logic (2026), First View, 1--13.
\href{https://doi.org/10.1017/jsl.2026.10200}
     {doi:10.1017/jsl.2026.10200}.

\bibitem{RYY}
Assaf Rinot, Zhixing You, and Jiachen Yuan,
\emph{A new model for all \(C\)-sequences are trivial},
The Journal of Symbolic Logic (2026), First View, 1--19.
\href{https://doi.org/10.1017/jsl.2026.10183}
     {doi:10.1017/jsl.2026.10183}.
Preprint (2025): \url{https://arxiv.org/abs/2504.06794}.

\bibitem{Schuerz}
Johannes Philipp Sch\"urz,
\emph{The cofinality of the strong measure zero ideal for
                   \(\kappa\) inaccessible},
Mathematical Logic Quarterly \textbf{69}(1) (2023), 31--39.
\href{https://doi.org/10.1002/malq.202000093}
     {doi:10.1002/malq.202000093}.

\end{thebibliography}

\end{document}
